\documentclass[11pt,a4paper]{article}
\usepackage[T1]{fontenc}
\usepackage[utf8]{inputenc}
\usepackage[brazilian]{babel}
\usepackage{mathtools,amsthm}
\usepackage{newpxtext,newpxmath}
\usepackage{booktabs,array,microtype,siunitx}
\usepackage[dvipsnames]{xcolor}
\usepackage[a4paper,margin=25mm,headheight=15pt,headsep=8mm,footskip=11mm]{geometry}
\usepackage{tikz,pgfplots,pgfplotstable}
\usetikzlibrary{arrows.meta,calc,positioning}
\pgfplotsset{compat=1.18,/pgf/number format/use comma}
\usepackage{caption,subcaption,fancyhdr,enumitem,needspace,float}
\usepackage{titlesec,setspace}
\usepackage[most]{tcolorbox}
\usepackage{algorithm,algpseudocode}
\usepackage[hidelinks]{hyperref}
\usepackage[nameinlink,noabbrev]{cleveref}
% Âncoras distintas para as linhas de cada algoritmo.
\makeatletter\newcounter{algorithmicH}\let\oldalgorithmic\algorithmic
\renewcommand{\algorithmic}{\stepcounter{algorithmicH}\oldalgorithmic}
\renewcommand{\theHALG@line}{ALG@line.\thealgorithmicH.\arabic{ALG@line}}\makeatother
\hypersetup{colorlinks=true,linkcolor=MidnightBlue,urlcolor=MidnightBlue,citecolor=MidnightBlue,
  pdftitle={O solver do TSPN: ciclo convexo, branch-and-bound e busca local},
  pdfauthor={Gabriel Freire Ushijima}}
\sisetup{output-decimal-marker={,},group-separator={.},group-minimum-digits=4}

% Paleta compartilhada com o relatório do oráculo convexo.
\definecolor{ink}{HTML}{223247}
\definecolor{navy}{HTML}{274568}
\definecolor{pOne}{HTML}{3677AD}
\definecolor{pTwo}{HTML}{C78A24}
\definecolor{pThree}{HTML}{6A4C9C}
\definecolor{good}{HTML}{007E72}
\definecolor{bad}{HTML}{B53E62}
\definecolor{quiet}{HTML}{617384}
\color{ink}
\tikzset{>=Latex,polyA/.style={draw=pOne,fill=pOne!9,line width=.9pt},
  polyB/.style={draw=pTwo,fill=pTwo!10,line width=.9pt},
  route/.style={draw=good,line width=1.4pt},wrong/.style={draw=bad,dashed,line width=1pt},
  aux/.style={draw=quiet,dashed,line width=.8pt},
  smalllabel/.style={font=\scriptsize,fill=white,fill opacity=.9,text opacity=1,inner sep=1pt}}

\setlength{\parindent}{0pt}\setlength{\parskip}{6.5pt}
\setlength{\emergencystretch}{2em}
\setstretch{1.07}
\captionsetup{font=small,labelfont={bf,color=navy},skip=6pt}
\titleformat{\section}{\Large\bfseries\color{navy}}{\thesection}{.7em}{}
\titleformat{\subsection}{\large\bfseries\color{good}}{\thesubsection}{.6em}{}
\titleformat{\paragraph}[runin]{\bfseries\color{navy}}{}{}{}[.]
\titlespacing*{\section}{0pt}{20pt plus 4pt minus 3pt}{8pt}
\titlespacing*{\subsection}{0pt}{14pt plus 4pt minus 2pt}{5pt}
\titlespacing*{\paragraph}{0pt}{5pt plus 2pt}{.6em}
\pagestyle{fancy}\fancyhf{}
\fancyhead[L]{\small\sffamily\color{navy}O SOLVER DO TSPN}
\fancyhead[R]{\small\sffamily\color{quiet}TouringPolygons}
\fancyfoot[L]{\small\color{quiet}Relatório técnico, 10 de outubro de 2026}
\fancyfoot[R]{\small\color{navy}\thepage}
\renewcommand{\headrulewidth}{.4pt}

\newtcolorbox{ideia}[1][]{colback=good!5,colframe=good!70!black,boxrule=.6pt,arc=2pt,
  left=6pt,right=6pt,top=4pt,bottom=4pt,fonttitle=\bfseries\small,title={#1}}
\newtcolorbox{fonte}{colback=navy!4,colframe=navy!35,boxrule=.4pt,arc=2pt,
  left=6pt,right=6pt,top=3pt,bottom=3pt,fontupper=\small}
\newtcolorbox{cuidado}[1][]{colback=bad!4,colframe=bad!60,boxrule=.6pt,arc=2pt,
  left=6pt,right=6pt,top=4pt,bottom=4pt,fonttitle=\bfseries\small,title={#1}}

\theoremstyle{plain}\newtheorem{proposicao}{Proposição}
\theoremstyle{definition}\newtheorem{definicao}{Definição}
\theoremstyle{plain}\newtheorem{teorema}{Teorema}
\theoremstyle{plain}\newtheorem{lema}{Lema}
\crefname{proposicao}{Proposição}{Proposições}
\crefname{teorema}{Teorema}{Teoremas}
\crefname{lema}{Lema}{Lemas}
\crefname{definicao}{Definição}{Definições}
\crefname{algorithm}{Algoritmo}{Algoritmos}
\crefname{figure}{Figura}{Figuras}
\crefname{table}{Tabela}{Tabelas}
\crefname{section}{Seção}{Seções}
\crefname{subsection}{Seção}{Seções}
\crefname{appendix}{Apêndice}{Apêndices}
\crefname{equation}{Equação}{Equações}
\creflabelformat{equation}{#2(#1)#3}
\AtBeginDocument{\renewcommand{\crefpairconjunction}{ e }\renewcommand{\crefmiddleconjunction}{, }%
  \renewcommand{\creflastconjunction}{ e }\renewcommand{\crefpairgroupconjunction}{ e }%
  \renewcommand{\crefmiddlegroupconjunction}{, }\renewcommand{\creflastgroupconjunction}{ e }}

\floatname{algorithm}{Algoritmo}
\algrenewcommand\algorithmicif{\textbf{se}}
\algrenewcommand\algorithmicthen{\textbf{então}}
\algrenewcommand\algorithmicelse{\textbf{senão}}
\algrenewcommand\algorithmicend{\textbf{fim}}
\algrenewcommand\algorithmicfor{\textbf{para}}
\algrenewcommand\algorithmicdo{\textbf{faça}}
\algrenewcommand\algorithmicwhile{\textbf{enquanto}}
\algrenewcommand\algorithmicreturn{\textbf{devolva}}
\algrenewcommand\algorithmicfunction{\textbf{função}}
\algrenewcommand\algorithmicrequire{\textbf{Entrada:}}
\algrenewcommand\algorithmicensure{\textbf{Saída:}}
\algrenewtext{EndIf}{\textbf{fim se}}
\algrenewtext{EndFor}{\textbf{fim para}}
\algrenewtext{EndWhile}{\textbf{fim enquanto}}
\algrenewtext{EndFunction}{\textbf{fim função}}
\algrenewcommand\algorithmicloop{\textbf{repita}}
\algrenewtext{EndLoop}{\textbf{fim repita}}
\algnewcommand{\IfThen}[2]{\State \algorithmicif\ #1\ \algorithmicthen\ #2}
\algnewcommand{\ForOne}[2]{\State \algorithmicfor\ #1\ \algorithmicdo\ #2}
\algrenewcommand{\algorithmiccomment}[1]{\hfill{\color{quiet}\(\triangleright\) #1}}

\newcommand{\R}{\mathbb R}
\newcommand{\norm}[1]{\lVert#1\rVert}
\newcommand{\abs}[1]{\lvert#1\rvert}
\newcommand{\OPT}{\mathrm{OPT}}
\newcommand{\code}[1]{\texttt{\small #1}}
\newcommand{\fl}{\mathrm{fl}}
\newcommand{\supp}{\mathrm{S}}
\DeclareMathOperator{\dist}{dist}

% Figuras do relatório. Desenhos geométricos usam a mesma escala nos dois
% eixos e as coordenadas exatas dos exemplos do texto (conferidas com o solver);
% os esquemas estão marcados como tais nas legendas.
\tikzset{coordgrid/.style={draw=quiet!22,line width=.25pt},
  polyC/.style={draw=pThree,fill=pThree!9,line width=.9pt},
  polyD/.style={draw=good,fill=good!9,line width=.9pt},
  dualvec/.style={draw=pThree,line width=1.1pt,-{Latex[length=2.4mm]}},
  coefvec/.style={draw=bad,line width=1.1pt,-{Latex[length=2.4mm]}},
  parent/.style={draw=quiet,dashed,line width=1pt},
  box/.style={draw=navy!70,fill=white,rounded corners=2pt,align=center,
    font=\small,inner sep=5pt,minimum height=9mm},
  okbox/.style={draw=good,fill=good!9,rounded corners=2pt,align=center,
    font=\small,inner sep=4pt},
  flow/.style={-{Latex[length=2.2mm]},draw=ink,line width=.7pt},
  region/.style={draw=pOne,fill=pOne!9,line width=.8pt},
  contact/.style={fill=ink},
  gapmark/.style={draw=bad,line width=1.2pt}}

% ---------------------------------------------------------------------------
% Visão geral (esquema).
\newcommand{\figvisaogeral}{%
\begin{tikzpicture}[node distance=6mm and 7mm]
\node[box,minimum width=30mm] (ils) {\textbf{ILS} (opcional)\\[1pt]\footnotesize ordem + um contato\\\footnotesize por região};
\node[box,right=9mm of ils,minimum width=34mm] (inc) {incumbente \(U^\star\)\\[1pt]\footnotesize tour validado};
\node[box,below=8mm of ils,minimum width=30mm] (fr) {fronteira\\[1pt]\footnotesize sequências parciais\\\footnotesize por limite inferior};
\node[box,right=9mm of fr,minimum width=34mm] (ex) {expansão do nó\\[1pt]\footnotesize região mais distante:\\\footnotesize inserir ou refinar};
\node[box,right=9mm of ex,minimum width=34mm] (tr) {\textbf{triagem}\\[1pt]\footnotesize limites de inserção\\\footnotesize (dual do pai)};
\node[box,below=8mm of tr,minimum width=34mm,draw=good,line width=.9pt] (or) {\textbf{oráculo de ciclo}\\[1pt]\footnotesize ordem fixa, regiões convexas\\\footnotesize \(L\le\mathrm{OPT}\le U\)};
\draw[flow] (ils)--(inc);
\draw[flow] (fr)--(ex);
\draw[flow] (ex)--(tr);
\draw[flow] (tr)--node[right,font=\scriptsize,text=quiet]{filhos não podados}(or);
\draw[flow] (or.west)-|node[pos=.25,below,font=\scriptsize,text=quiet]{limite \(L\) do filho}(fr.south);
\draw[flow] (or.north east)--++(4mm,0)|-node[pos=.75,above,font=\scriptsize,text=quiet]{tour que cobre tudo}(inc.east);
\draw[flow,dashed] (inc.south)--node[right,font=\scriptsize,text=quiet]{corte \(U^\star-g\)}(ex.north);
\end{tikzpicture}}

% ---------------------------------------------------------------------------
% Exemplo das três caixas: ciclo ótimo e coeficientes do certificado dual.
\newcommand{\figtrescaixas}{%
\begin{tikzpicture}[scale=.62]
\begin{scope}
  \draw[coordgrid] (-3.5,-1.5) grid (3.5,4.5);
  \draw[polyA] (-3,-1) rectangle (-2,1);
  \draw[polyB] (2,-1) rectangle (3,1);
  \draw[polyC] (-1,3) rectangle (1,4);
  \draw[route,-{Latex[length=2.4mm]}] (-2,1)--(2,1);
  \draw[route,-{Latex[length=2.4mm]}] (2,1)--(0,3);
  \draw[route,-{Latex[length=2.4mm]}] (0,3)--(-2,1);
  \foreach \p in {(-2,1),(2,1),(0,3)} \fill[contact] \p circle (3pt);
  \node[smalllabel] at (-2.5,-.3) {\(P_1\)};
  \node[smalllabel] at (2.5,-.3) {\(P_2\)};
  \node[smalllabel] at (0,3.6) {\(P_3\)};
  \node[smalllabel,above left] at (-2,1.05) {\(p_1\)};
  \node[smalllabel,above right] at (2,1.05) {\(p_2\)};
  \node[smalllabel,below] at (0,2.9) {\(p_3\)};
  \node[smalllabel,below] at (0,.95) {\(u_1\)};
  \node[smalllabel,right] at (1.05,2.05) {\(u_2\)};
  \node[smalllabel,left] at (-1.05,2.05) {\(u_3\)};
  \node[font=\scriptsize\bfseries,color=navy] at (0,-2.2) {(a) ciclo ótimo, comprimento \(4+4\sqrt2\)};
\end{scope}
\begin{scope}[xshift=8.6cm]
  \draw[coordgrid] (-3.5,-1.5) grid (3.5,4.5);
  \draw[polyA] (-3,-1) rectangle (-2,1);
  \draw[polyB] (2,-1) rectangle (3,1);
  \draw[polyC] (-1,3) rectangle (1,4);
  \draw[parent] (-2,1)--(2,1)--(0,3)--cycle;
  \draw[coefvec] (-2,1)--++(-1.2071,-0.5);
  \draw[coefvec] (2,1)--++(1.2071,-0.5);
  \draw[coefvec] (0,3)--++(0,1);
  \foreach \p in {(-2,1),(2,1),(0,3)} \fill[contact] \p circle (3pt);
  \node[smalllabel,left] at (-3.2,.5) {\(c_1\)};
  \node[smalllabel,right] at (3.2,.5) {\(c_2\)};
  \node[smalllabel,right] at (.08,4.1) {\(c_3\)};
  \node[font=\scriptsize\bfseries,color=navy] at (0,-2.2) {(b) \(c_i=u_{i-1}-u_i\) (escala \(1/\sqrt2\))};
\end{scope}
\end{tikzpicture}}

% ---------------------------------------------------------------------------
% Função âncora sobre a aresta inferior de P_3.
\newcommand{\figancora}{%
\begin{tabular}{@{}c@{\hspace{9mm}}c@{}}
\begin{tikzpicture}[scale=.62,baseline=(current bounding box.south)]
  \draw[coordgrid] (-3.5,-1.5) grid (3.5,4.5);
  \draw[polyA] (-3,-1) rectangle (-2,1);
  \draw[polyB] (2,-1) rectangle (3,1);
  \draw[polyC] (-1,3) rectangle (1,4);
  \draw[draw=navy,line width=2.2pt] (-1,3)--(1,3);
  \draw[parent] (-2,1)--(2,1)--(0,3)--cycle;
  \draw[route,-{Latex[length=2.4mm]}] (0.6,3)--(-2,1);
  \draw[route,-{Latex[length=2.4mm]}] (-2,1)--(2,1);
  \draw[route,-{Latex[length=2.4mm]}] (2,1)--(0.6,3);
  \fill[bad] (0.6,3) circle (3.2pt);
  \node[smalllabel,right] at (1.1,3.15) {\(x=(t,3)\)};
  \node[smalllabel] at (-2.5,-.3) {\(P_1\)};
  \node[smalllabel] at (2.5,-.3) {\(P_2\)};
  \node[smalllabel] at (0,3.6) {\(P_3\)};
  \node[smalllabel,text=navy,below] at (-.55,2.92) {\(e\)};
\end{tikzpicture}
&
\begin{tikzpicture}[baseline=(current bounding box.south)]
\begin{axis}[width=74mm,height=52mm,
  xmin=-1.08,xmax=1.08,ymin=9.6,ymax=9.92,xlabel={\(t\)},ylabel={\(F(t)\)},
  xtick={-1,-0.5,0,0.5,1},ytick={9.6,9.7,9.8,9.9},
  yticklabel style={/pgf/number format/fixed,/pgf/number format/precision=1},
  tick label style={font=\scriptsize},label style={font=\small},
  axis lines=left,grid=major,grid style={quiet!18}]
  \addplot[navy,line width=1.1pt,domain=-1:1,samples=101] {sqrt((x+2)^2+4)+4+sqrt((2-x)^2+4)};
  \addplot[only marks,mark=*,mark size=1.8pt,good] coordinates {(0,9.656854)};
  \addplot[only marks,mark=*,mark size=1.8pt,bad] coordinates {(0.6,9.721555)};
  \addplot[pTwo,dashed,line width=.9pt,domain=-1:-0.7] {9.841619-0.384900*(x+1)};
  \addplot[pTwo,dashed,line width=.9pt,domain=0.7:1] {9.841619+0.384900*(x-1)};
  \node[font=\scriptsize,anchor=north] at (axis cs:0,9.651) {\(4+4\sqrt2\)};
  \node[font=\scriptsize,anchor=west] at (axis cs:-1.02,9.895) {\(f'(-1)<0\)};
  \node[font=\scriptsize,anchor=east] at (axis cs:1.02,9.895) {\(f'(1)>0\)};
\end{axis}
\end{tikzpicture}\\[2pt]
{\scriptsize\bfseries\color{navy}(a) caminho ancorado em \(x\), \(t=0{,}6\)} &
{\scriptsize\bfseries\color{navy}(b) \(F\) ao longo de \(e\)}
\end{tabular}}

% ---------------------------------------------------------------------------
% Posições de inserção: caminho com extremos fixos e ciclo (esquema).
\newcommand{\figlacunas}{%
\begin{tikzpicture}[x=9mm,y=9mm]
\begin{scope}
  \fill[ink] (0,0) circle (2.4pt) node[below=2pt,font=\scriptsize] {\(s\)};
  \fill[ink] (8,0) circle (2.4pt) node[below=2pt,font=\scriptsize] {\(t\)};
  \foreach \x/\l in {2/1,4/2,6/3}{\draw[region] (\x,0) circle (.42);\node[font=\scriptsize] at (\x,0) {\(P_{\l}\)};}
  \foreach \a/\b in {0/1.58,2.42/3.58,4.42/5.58,6.42/8}{\draw[route] (\a,0)--(\b,0);}
  \foreach \x/\l in {1/0,3/1,5/2,7/3}{\draw[gapmark] (\x,-.22)--(\x,.22);\node[font=\scriptsize,text=bad] at (\x,.5) {\(\l\)};}
  \node[font=\scriptsize\bfseries,color=navy] at (4,-1.3) {(a) caminho: \(m+1=4\) posições};
\end{scope}
\begin{scope}[xshift=10.5cm]
  \foreach \ang/\l in {90/1,210/2,330/3}{\draw[region] (\ang:1.6) circle (.42);\node[font=\scriptsize] at (\ang:1.6) {\(P_{\l}\)};}
  \draw[route] (90:1.6) ++(250:.42) -- ($(210:1.6)+(50:.42)$);
  \draw[route] (210:1.6) ++(-10:.42) -- ($(330:1.6)+(190:.42)$);
  \draw[route] (330:1.6) ++(130:.42) -- ($(90:1.6)+(-70:.42)$);
  \foreach \ang/\l in {150/1,270/2,30/3}{\draw[gapmark] (\ang:.68)--(\ang:1.12);\node[font=\scriptsize,text=bad] at (\ang:1.38) {\(\l\)};}
  \node[font=\scriptsize\bfseries,color=navy] at (0,-1.98) {(b) ciclo: \(m=3\) posições};
\end{scope}
\end{tikzpicture}}

% ---------------------------------------------------------------------------
% Limite de inserção: o pai (três caixas) e a caixa R inserida na lacuna 1.
\newcommand{\figinsercao}{%
\begin{tikzpicture}[scale=.62]
\begin{scope}
  \draw[coordgrid] (-3.5,-3.5) grid (3.5,4.5);
  \draw[polyA] (-3,-1) rectangle (-2,1);
  \draw[polyB] (2,-1) rectangle (3,1);
  \draw[polyC] (-1,3) rectangle (1,4);
  \draw[polyD] (-0.5,-3) rectangle (0.5,-2);
  \draw[parent] (-2,1)--(2,1)--(0,3)--cycle;
  \draw[dualvec] (-2,1)--(0,-2);
  \draw[dualvec] (0,-2)--(2,1);
  \foreach \p in {(-2,1),(2,1),(0,3)} \fill[contact] \p circle (3pt);
  \fill[bad] (0,-2) circle (3.2pt);
  \node[smalllabel] at (-2.5,-.3) {\(P_1\)};
  \node[smalllabel] at (2.5,-.3) {\(P_2\)};
  \node[smalllabel] at (0,3.6) {\(P_3\)};
  \node[smalllabel] at (0,-2.6) {\(R\)};
  \node[smalllabel,left] at (-2.05,1.15) {\(q_1\)};
  \node[smalllabel,right] at (2.05,1.15) {\(q_2\)};
  \node[smalllabel,right] at (.12,-1.75) {\(c\)};
  \node[smalllabel,left] at (-1.05,-.5) {\(a\)};
  \node[smalllabel,right] at (1.05,-.5) {\(b\)};
  \node[font=\scriptsize\bfseries,color=navy] at (0,-4.2) {(a) dual do pai e contato proposto};
\end{scope}
\begin{scope}[xshift=8.6cm]
  \draw[coordgrid] (-3.5,-3.5) grid (3.5,4.5);
  \draw[polyA] (-3,-1) rectangle (-2,1);
  \draw[polyB] (2,-1) rectangle (3,1);
  \draw[polyC] (-1,3) rectangle (1,4);
  \draw[polyD] (-0.5,-3) rectangle (0.5,-2);
  \draw[route,-{Latex[length=2.4mm]}] (-2,0.5)--(0,-2);
  \draw[route,-{Latex[length=2.4mm]}] (0,-2)--(2,0.5);
  \draw[route,-{Latex[length=2.4mm]}] (2,0.5)--(0,3);
  \draw[route,-{Latex[length=2.4mm]}] (0,3)--(-2,0.5);
  \foreach \p in {(-2,0.5),(0,-2),(2,0.5),(0,3)} \fill[contact] \p circle (3pt);
  \node[font=\scriptsize\bfseries,color=navy] at (0,-4.2) {(b) ótimo do filho, \(2\sqrt{41}\)};
\end{scope}
\end{tikzpicture}}

% ---------------------------------------------------------------------------
% Piso da fronteira: o limite inferior congela quando a memória acaba.
\newcommand{\figpiso}{%
\begin{tikzpicture}
\begin{axis}[width=118mm,height=58mm,xmin=0,xmax=300,ymin=3.4,ymax=8,
  xlabel={tempo de B\&B (s)},ylabel={gap certificado (\%)},
  xtick={0,60,120,180,240,300},tick label style={font=\scriptsize},label style={font=\small},
  axis lines=left,grid=major,grid style={quiet!18},
  legend style={font=\scriptsize,draw=none,fill=white,fill opacity=.9,text opacity=1,at={(.98,.96)},anchor=north east}]
  \addplot[bad,line width=1.1pt,mark=*,mark size=1.2pt] table[x=t,y=small]{\dadospiso};
  \addplot[good,line width=1.1pt,mark=*,mark size=1.2pt] table[x=t,y=large]{\dadospiso};
  \legend{fronteira de 256 MiB,fronteira de 6 GiB}
\end{axis}
\end{tikzpicture}}

% ---------------------------------------------------------------------------
% O subproblema de um contato: a mesma fórmula para reflexão e cruzamento.
\newcommand{\figmelhorcontato}{%
\begin{tikzpicture}[scale=.78]
\begin{scope}
  \draw[coordgrid] (-.5,-.5) grid (6.5,3.5);
  \draw[quiet,dashed] (-.5,2)--(6.5,2);
  \draw[draw=navy,line width=2.2pt] (1,2)--(5,2);
  \draw[route,-{Latex[length=2.4mm]}] (0,0)--(4,2);
  \draw[route,-{Latex[length=2.4mm]}] (4,2)--(6,1);
  \draw[wrong] (4,2)--(6,3);
  \draw[aux] (0,0)--(0,2);\draw[aux] (6,1)--(6,2);
  \fill[ink] (0,0) circle (2.4pt);\fill[ink] (6,1) circle (2.4pt);\fill[quiet] (6,3) circle (2pt);
  \fill[bad] (4,2) circle (2.8pt);
  \node[smalllabel,below] at (0,-.05) {\(\ell\)};
  \node[smalllabel,below] at (6,.95) {\(r\)};
  \node[smalllabel,right] at (6.05,3) {\(r'\)};
  \node[smalllabel,above] at (4,2.08) {\(p=(4,2)\)};
  \node[smalllabel,left] at (-.05,1) {\(h_\ell=2\)};
  \node[smalllabel,right] at (6.05,1.5) {\(h_r=1\)};
  \node[smalllabel,text=navy] at (2,2.3) {aresta};
  \node[font=\scriptsize\bfseries,color=navy] at (3,-1.2) {(a) mesmo lado: reflexão};
\end{scope}
\begin{scope}[xshift=8.2cm]
  \draw[coordgrid] (-.5,-.5) grid (6.5,3.5);
  \draw[quiet,dashed] (-.5,2)--(6.5,2);
  \draw[draw=navy,line width=2.2pt] (1,2)--(5,2);
  \draw[route,-{Latex[length=2.4mm]}] (0,0)--(4,2);
  \draw[route,-{Latex[length=2.4mm]}] (4,2)--(6,3);
  \draw[aux] (0,0)--(0,2);\draw[aux] (6,3)--(6,2);
  \fill[ink] (0,0) circle (2.4pt);\fill[ink] (6,3) circle (2.4pt);
  \fill[bad] (4,2) circle (2.8pt);
  \node[smalllabel,below] at (0,-.05) {\(\ell\)};
  \node[smalllabel,right] at (6.05,3) {\(r\)};
  \node[smalllabel,below] at (4,1.92) {\(p=(4,2)\)};
  \node[smalllabel,left] at (-.05,1) {\(h_\ell=2\)};
  \node[smalllabel,right] at (6.05,2.5) {\(h_r=1\)};
  \node[font=\scriptsize\bfseries,color=navy] at (3,-1.2) {(b) lados opostos: cruzamento};
\end{scope}
\end{tikzpicture}}

% ---------------------------------------------------------------------------
% A armadilha da descida por coordenadas: contatos coincidentes.
\newcommand{\figarmadilha}{%
\begin{tikzpicture}[scale=.82]
  \draw[coordgrid] (-.5,-.8) grid (6.5,3.8);
  \draw[polyA,fill opacity=.85] (1,1) rectangle (4,3);
  \draw[polyB,fill opacity=.6] (2,1) rectangle (5,3);
  \draw[wrong,-{Latex[length=2.2mm]}] (0,0)--(4,3);
  \draw[wrong,-{Latex[length=2.2mm]}] (4,3)--(6,0);
  \draw[route,-{Latex[length=2.4mm]}] (0,0)--(3,1);
  \draw[route,-{Latex[length=2.4mm]}] (3,1)--(6,0);
  \fill[ink] (0,0) circle (2.4pt);\fill[ink] (6,0) circle (2.4pt);
  \fill[bad] (4,3) circle (2.8pt);\fill[good] (3,1) circle (2.8pt);
  \node[smalllabel,below left] at (0,0) {\(a\)};
  \node[smalllabel,below right] at (6,0) {\(b\)};
  \node[smalllabel,above] at (4,3.08) {\(z=(4,3)\)};
  \node[smalllabel,below] at (3,.9) {\((3,1)\)};
  \node[smalllabel] at (1.5,2.6) {\(P\)};
  \node[smalllabel] at (4.6,2.6) {\(Q\)};
\end{tikzpicture}}

% ---------------------------------------------------------------------------
% Reotimização exata por janela (esquema).
\newcommand{\figjanela}{%
\begin{tikzpicture}
  \def\n{16}
  \foreach \i in {0,...,15}{
    \pgfmathsetmacro{\ang}{90-360*\i/\n}
    \coordinate (c\i) at (\ang:2.2);}
  \foreach \i [evaluate=\i as \j using {int(mod(\i+1,16))}] in {0,...,15}{\draw[route] (c\i)--(c\j);}
  \foreach \i in {2,...,9}{\draw[region,fill=pTwo!18,draw=pTwo] (c\i) circle (.2);}
  \foreach \i in {0,1,10,11,12,13,14,15}{\draw[region] (c\i) circle (.2);}
  \fill[ink] (c1) circle (2.4pt);\fill[ink] (c10) circle (2.4pt);
  \node[font=\scriptsize,anchor=south] at ($(c1)+(0,.25)$) {contato fixo};
  \node[font=\scriptsize,anchor=north] at ($(c10)+(0,-.25)$) {contato fixo};
  \node[font=\scriptsize,align=center,text=pTwo!80!black] at (2.3,-2.8) {janela de 8 posições:\\ B\&B de ordem livre entre os fixos};
  \node[font=\scriptsize,align=center,text=quiet] at (0,0) {resto do tour\\ inalterado};
\end{tikzpicture}}

% Dados das figuras de medição. Origem de cada tabela no comentário acima dela;
% formulação, protocolo e limitações das campanhas em
% benchmarks/results-saved/README.md.

% Piso da fronteira (10/10/2026, Mac M4 Pro, binário 0722d66): 100i1000-410
% (coleção da Paula), B&B de 300 s com --cycle-optimization lazy --lazy-oracles
% --dive-interval 0, a partir do mesmo tour (5.427,10), gap relativo 1e-6 e
% cache,features,root,interval; as duas fronteiras rodaram juntas, um
% processo cada. Gap = 1 - LB/UB em %, lido dos relatórios de progresso de 15 s.
\pgfplotstableread{
t     small  large
15    7.56   7.48
30    6.55   6.46
45    6.00   5.89
60    5.63   5.50
75    5.35   5.20
90    5.12   4.96
105   4.92   4.76
120   4.88   4.58
135   4.88   4.43
150   4.88   4.29
165   4.88   4.17
180   4.88   4.03
195   4.88   3.96
210   4.88   3.90
225   4.88   3.85
240   4.88   3.79
255   4.88   3.74
270   4.88   3.69
285   4.88   3.64
}\dadospiso


\begin{document}
\pagenumbering{Alph}

\begin{titlepage}
\vspace*{18mm}
{\color{quiet}\sffamily\small TouringPolygons --- relatório técnico\par}
\vspace{6mm}
{\color{navy}\bfseries\fontsize{24}{29}\selectfont O solver do TSPN\par}
\vspace{3mm}
{\color{good}\Large O ciclo convexo de ordem fixa, o branch-and-bound de inserções e a busca local iterada\par}
\vspace{10mm}
{\large Gabriel Freire Ushijima\par}
{\color{quiet}10 de outubro de 2026\par}
\vspace{14mm}
\begin{tcolorbox}[colback=navy!3,colframe=navy!40,boxrule=.5pt,arc=2pt,left=8pt,right=8pt,top=6pt,bottom=6pt]
\textbf{Resumo.} O TSPN pede o tour fechado mais curto que visita regiões do plano (pontos, segmentos e polígonos simples, possivelmente não convexos e sobrepostos) em qualquer ordem e sem ponto de partida. Nosso solver é o mesmo \emph{branch-and-bound} do TPP de ordem livre, com quatro mudanças: a topologia (ciclo em vez de caminho), o subproblema de cada nó (um ciclo de ordem fixa por regiões convexas, que não tem origem), a triagem dos filhos (limites de inserção tirados do dual do pai e provados em \texttt{binary64}) e uma busca local iterada que alimenta o incumbente. Este relatório explica as três peças supondo conhecido o TPP de ordem livre: como o ciclo de ordem fixa se reduz a caminhos ancorados na fronteira de um polígono; como a certificação primal-dual em ponto flutuante se adapta ao ciclo; por que um ponto na instância transforma o TSPN num TPP; como o dual do pai limita os filhos sem resolvê-los; e como a busca local combina vizinhanças clássicas com reotimização exata por janelas. Nas medições, a regressão de 48 casos SoCG caiu de 131 para 2,7 s, e nas 235 instâncias da Paula temos 220 ótimos certificados, nunca com um tour pior que o melhor dela; nas 15 restantes, o limite inferior é limitado pela memória da fronteira, que reduzimos em seguida.
\end{tcolorbox}
\vfill
{\small\color{quiet}Código descrito: \code{main} em \code{0722d66}, mais a campanha \code{paula-campaign} e a fronteira compacta de 12/10 (branches \code{paula-long-campaign} e \code{frontier-compact}). Este documento pressupõe o relatório \emph{O oráculo convexo certificado}~\cite{oraculo} e o algoritmo de ordem livre~\cite{ordemlivre}; o que já está lá é só resumido.\par}
\end{titlepage}

\pagenumbering{roman}
\setcounter{tocdepth}{2}
{\small\setlength{\parskip}{0pt}\tableofcontents}
\clearpage
\pagenumbering{arabic}

%==========================================================================
\section{Introdução}

\subsection{O problema}
Recebemos regiões fechadas \(R_1,\dots,R_n\subset\R^2\). Cada uma é um ponto (um vértice), um segmento fechado (dois vértices) ou um polígono simples de área positiva, convexo ou não; as regiões podem se sobrepor, se tocar ou conter umas às outras. Um \emph{tour} é uma poligonal fechada que passa por todas as regiões; queremos o mais curto. Não há depósito nem ordem prescrita, e o tour pode se cruzar. Com \(n\le1\) o ótimo é zero.

É o mesmo problema que Fekete et al.~\cite{fekete} resolvem por \emph{branch-and-bound}, e o mesmo do rascunho da Paula~\cite{paula}, que o ataca com uma busca local iterada (ILS-BCD) e um modelo MIQCP para o CPLEX. A coleção dela, 235 instâncias tiradas da GTSP-Lib e da MOM-Lib, é a nossa principal referência de avaliação.

\subsection{O que muda em relação ao TPP de ordem livre}
O leitor já conhece o solver do TPP de ordem livre com extremos fixos~\cite{ordemlivre,oraculo}: um nó guarda uma sequência parcial de regiões convexas (fechos ou peças da decomposição); o oráculo resolve o caminho de \(s\) a \(t\) por essa sequência; a região original mais distante do caminho relaxado é inserida em todas as posições ou refinada em peças. O TSPN usa \emph{essa mesma busca}: ela tem uma opção de topologia, e não existe um segundo \emph{branch-and-bound}. A \cref{tab:diferencas} resume o que muda.

\begin{table}[ht]\centering\small
\begin{tabular}{@{}>{\raggedright\arraybackslash}p{32mm}>{\raggedright\arraybackslash}p{50mm}>{\raggedright\arraybackslash}p{62mm}@{}}\toprule
& TPP de ordem livre & TSPN\\\midrule
Extremos & \(s,t\) fixos & nenhum; o tour é fechado\\
Raiz & sequência vazia, caminho reto \(s\to t\) & uma região (ou três, com \code{root}); \cref{sec:raiz}\\
Posições de inserção & \(m+1\) para \(m\) regiões & \(m\), incluindo o elo de fechamento; com \(m=2\), só uma\\
Subproblema do nó & caminho de ordem fixa: mapas de Dror com origem \(s\) & ciclo de ordem fixa, sem origem: \cref{sec:ciclo}\\
Limite dual & \((t-s)\cdot u_m+\sum_i\supp_i\) & só os termos de suporte, invariante por translação\\
Triagem dos filhos & limite de inserção do caminho & limite de inserção do ciclo, em \texttt{binary64}: \cref{sec:insercao}\\
Instância com um ponto & --- & o tour passa pelo ponto: vira TPP de \(p\) a \(p\) (\cref{sec:ancora})\\
Limite superior & guloso, 2-opt, varredura de contatos & os mesmos e, opcionalmente, a ILS (\cref{sec:ils})\\\bottomrule
\end{tabular}
\caption{O que o TSPN reaproveita e o que muda. Fila, mergulhos, decomposição preguiçosa, escolha da região, verificação de visitas e contabilidade da fronteira são compartilhados.}\label{tab:diferencas}
\end{table}

\subsection{Visão geral e roteiro}
A \cref{fig:visaogeral} mostra as peças. A \cref{sec:ciclo} trata do oráculo de ciclo: a formulação, o certificado dual, a redução à fronteira que substitui a origem que falta, a construção rápida que fecha quase todas as chamadas e o estágio \texttt{binary64}. A \cref{sec:bnb} trata da busca: âncora de ponto, raiz, ramificação, limites de inserção, política de busca e memória da fronteira. A \cref{sec:ils} trata da ILS. As \cref{sec:medicoes,sec:limites} trazem medições e limitações.

\begin{figure}[ht]\centering\figvisaogeral
\caption{Esquema do solver. A ILS só fornece incumbentes; a triagem poda filhos antes do oráculo; o oráculo devolve o limite do filho e um tour relaxado, que vira incumbente se visitar todas as regiões originais.}\label{fig:visaogeral}
\end{figure}

%==========================================================================
\section{O ciclo de ordem fixa por regiões convexas}\label{sec:ciclo}

\subsection{Formulação}
Sejam \(P_1,\dots,P_k\) (\(k\ge2\)) regiões convexas fechadas: pontos, segmentos ou polígonos convexos de área positiva. Um ciclo viável escolhe um contato \(p_i\in P_i\) por região; os índices são cíclicos, com \(p_{k+1}=p_1\). O problema é
\begin{equation}\label{eq:ciclo}
\OPT=\min_{p_1\in P_1,\dots,p_k\in P_k}\;\sum_{i=1}^{k}\norm{p_{i+1}-p_i}.
\end{equation}
Como no caminho, a função objetivo é convexa e o domínio é um produto de convexos compactos: o mínimo existe e todo mínimo local é global. Algumas propriedades que o caminho não tem, ou tem de outro jeito:
\begin{itemize}[leftmargin=*,itemsep=1pt]
\item \(\OPT=0\) se e só se \(P_1\cap\dots\cap P_k\ne\emptyset\): um ciclo de comprimento zero é um único ponto comum.
\item Rotacionar ou inverter a ordem não muda o problema; não há um ``primeiro'' contato.
\item Com \(k=2\), \(\OPT=2\dist(P_1,P_2)\).
\item Os contatos ótimos não precisam ser únicos. No exemplo da \cref{sec:insercao}, o ótimo do filho forma uma família de paralelogramos de mesmo comprimento.
\item O valor ótimo é contínuo nos dados: mover cada região no máximo \(\varepsilon\) (em distância de Hausdorff) muda \(\OPT\) no máximo \(2k\varepsilon\).
\end{itemize}

\paragraph{Exemplo} Sejam \(P_1=[-3,-2]\times[-1,1]\), \(P_2=[2,3]\times[-1,1]\) e \(P_3=[-1,1]\times[3,4]\), nessa ordem (\cref{fig:trescaixas}). O elo \(P_1\to P_2\) mede pelo menos 4, e fica com 4 se os dois contatos tiverem a mesma ordenada; os outros dois elos encurtam quando os contatos sobem. O ótimo usa os cantos \(p_1=(-2,1)\) e \(p_2=(2,1)\) e o ponto \(p_3=(0,3)\) da aresta inferior de \(P_3\), com \(\OPT=4+2\sqrt8=4+4\sqrt2\approx9{,}6569\). O solver dá o mesmo valor.

\begin{figure}[ht]\centering\figtrescaixas
\caption{O exemplo das três caixas. (a) O ciclo ótimo e as direções \(u_i\) dos seus elos. (b) Os coeficientes \(c_i=u_{i-1}-u_i\) do certificado dual em cada contato. Cada um é uma combinação não negativa das normais internas das arestas que contêm o contato, por isso o mínimo de \(v\mapsto v\cdot c_i\) em \(P_i\) é atingido no próprio contato.}\label{fig:trescaixas}
\end{figure}

\subsection{O certificado dual do ciclo}\label{sec:dualciclo}
Para vetores \(u_1,\dots,u_k\) com \(\norm{u_i}\le1\), um por elo, a desigualdade \(\norm d\ge u\cdot d\) somada em volta do ciclo dá, para qualquer ciclo viável \(q\),
\[
\sum_{i=1}^k\norm{q_{i+1}-q_i}\;\ge\;\sum_{i=1}^k u_i\cdot(q_{i+1}-q_i)\;=\;\sum_{i=1}^k q_i\cdot(u_{i-1}-u_i),
\]
reagrupando por contato. Os coeficientes \(c_i=u_{i-1}-u_i\) somam exatamente zero, então podemos subtrair de todos os contatos um mesmo ponto de referência \(r\) sem mudar a soma.

\begin{definicao}[Limite dual do ciclo]\label{def:dualciclo}
Para \(\norm{u_i}\le1\) e qualquer \(r\in\R^2\),
\begin{equation}\label{eq:dualciclo}
D(u)=\sum_{i=1}^k\;\min_{v\in P_i}\,(v-r)\cdot(u_{i-1}-u_i)\;\le\;\OPT .
\end{equation}
\end{definicao}

Comparando com o caminho~\cite{oraculo}, sumiu o termo \((t-s)\cdot u_m\) dos extremos, e com ele a dependência de uma origem. O valor de \(D\) não depende de \(r\); o código usa a média das médias de vértices das regiões, para que os produtos escalares fiquem pequenos e o erro de arredondamento também. Cada mínimo é atingido num vértice, inclusive num segmento (dois vértices) e num ponto (um).

No exemplo, as direções do ciclo ótimo são \(u_1=(1,0)\), \(u_2=(-1,1)/\sqrt2\) e \(u_3=(-1,-1)/\sqrt2\). Com \(r=0\):
\[
\begin{array}{lll}
c_1=u_3-u_1=\bigl(-1-\tfrac1{\sqrt2},-\tfrac1{\sqrt2}\bigr), & \min_{P_1} v\cdot c_1 \text{ em } (-2,1): & 2+\tfrac1{\sqrt2},\\[3pt]
c_2=u_1-u_2=\bigl(1+\tfrac1{\sqrt2},-\tfrac1{\sqrt2}\bigr), & \min_{P_2} v\cdot c_2 \text{ em } (2,1): & 2+\tfrac1{\sqrt2},\\[3pt]
c_3=u_2-u_3=(0,\sqrt2), & \min_{P_3} v\cdot c_3 \text{ na aresta } y=3: & 3\sqrt2,
\end{array}
\]
e \(D(u)=4+\sqrt2+3\sqrt2=4+4\sqrt2\). O limite inferior alcança o comprimento: está provado que o ciclo é ótimo.

\paragraph{Quando o certificado é exato} O argumento da \cref{fig:trescaixas}(b) vale em geral: se cada \(c_i\) satisfaz \(c_i\cdot(v-p_i)\ge0\) para todo \(v\in P_i\), então \(D(u)\) é o comprimento do ciclo \(p\), que portanto é ótimo. Para elos de comprimento positivo, \(u_i\) é forçado: é a direção do elo. Num elo de comprimento zero, qualquer \(\norm{u_i}\le1\) é admissível, e o certificado exato decide se existe uma escolha que satisfaz todas as condições do bloco de contatos coincidentes, com a mesma propagação disco/cone do oráculo de caminho~\cite{oraculo,certificadociclo}. A condição é necessária e suficiente (é a condição de otimalidade convexa do problema), e o verificador \code{tpp\_convex\_verify\_cycle\_certificate} a testa em aritmética exata. No estágio \texttt{binary64} (\cref{sec:binary64}), não se testa otimalidade: só se avalia \(D(u)\) com arredondamento dirigido, para um \(u\) proposto.

\subsection{O obstáculo: não há origem}
O algoritmo de mapas de último passo de Dror et al.~\cite{dror} calcula caminhos mais curtos \emph{a partir de uma origem fixa}: \(F_0(q)=\norm{q-s}\) é o primeiro nível da recorrência. No ciclo, a origem natural seria um contato ótimo, que é justamente o que não sabemos. Precisamos então de duas coisas: um conjunto finito de ``lugares'' onde procurar a origem, garantidamente com um contato ótimo; e um jeito de procurar nele sem enumerar pontos. A primeira vem da proposição abaixo; a segunda, da convexidade da função âncora (\cref{sec:funcaoancora}).

\begin{proposicao}[Basta a fronteira]\label{prop:fronteira}
Se \(P_1\cap\dots\cap P_k=\emptyset\), todo ciclo ótimo tem algum contato na fronteira de alguma região. (Um ponto ou um segmento é a sua própria fronteira.)
\end{proposicao}
\begin{proof}
Suponha um ciclo ótimo com todos os contatos no interior das suas regiões. Agrupe os contatos consecutivos coincidentes em blocos maximais. Como o comprimento é positivo (as regiões não têm ponto comum), há pelo menos dois blocos, e cada bloco tem um elo de entrada e um de saída de comprimento positivo, com direções \(u_{\text{in}}\) e \(u_{\text{out}}\). Transladar o bloco inteiro por um vetor pequeno mantém a viabilidade, porque todos os contatos estão em interiores; o gradiente do comprimento nessa translação é \(u_{\text{in}}-u_{\text{out}}\), e a otimalidade exige que ele seja zero. Então todos os elos positivos têm a mesma direção \(u\), e a soma dos vetores dos elos é \((\sum_i\norm{d_i})\,u\ne0\). Mas a soma dos elos de um ciclo fechado é zero: contradição.
\end{proof}

A proposição fala de \emph{alguma} região, não de uma região escolhida. Uma região que contém outra pode não ter contato ótimo na sua fronteira, e escolher só ela estaria errado. Por isso o fluxo exato varre as fronteiras de todas as regiões. No caso de polígonos dois a dois disjuntos, basta o menor: um contato interior fica entre dois vizinhos exteriores, em linha reta, e pode deslizar até a fronteira sem mudar o comprimento.

\subsection{A função âncora}\label{sec:funcaoancora}
Divida cada aresta nos pontos em que ela cruza as fronteiras das outras regiões. Num trecho aberto dessa subdivisão, o conjunto de regiões que contêm o ponto não muda. Para \(x\) na fronteira de \(P_j\), seja \(F(x)\) o comprimento do caminho mais curto que sai de \(x\), visita \(P_{j+1},\dots,P_k,P_1,\dots,P_{j-1}\) nessa ordem e volta a \(x\). É um TPP de ordem fixa com \(s=t=x\), que o oráculo de caminho resolve. Pela \cref{prop:fronteira}, o mínimo de \(F\) sobre todas as fronteiras é \(\OPT\).

\begin{proposicao}[Convexidade da âncora]\label{prop:ancora}
Num trecho \(x=a+t\,e\) da subdivisão, \(f(t)=F(a+te)\) é convexa e 2-Lipschitz. No interior do trecho, \(f\) é derivável e
\begin{equation}\label{eq:derivada}
f'(t)=(u_{\text{in}}-u_{\text{out}})\cdot e,
\end{equation}
onde \(u_{\text{in}}\) é a direção do último elo positivo que chega a \(x\) e \(u_{\text{out}}\) a do primeiro que sai.
\end{proposicao}
\begin{proof}[Ideia da prova]
A soma de normas é conjuntamente convexa em \((x,p)\), e a minimização parcial em \(p\) preserva convexidade; mover \(x\) por \(\delta\) muda os dois elos incidentes em no máximo \(\norm\delta\) cada. As regiões que contêm \(x\) podem ser visitadas em \(x\) sem custo; elas formam um prefixo e um sufixo da ordem que não mudam no trecho, então os elos incidentes que ``contam'' ligam \(x\) a regiões que não o contêm e têm comprimento positivo. Duas soluções ótimas não podem ter direções diferentes nesses elos (igualdade na desigualdade triangular força direções iguais), o que dá a derivada única. A prova completa está em~\cite{ciclo}.
\end{proof}

\paragraph{Exemplo} Na aresta inferior \(e\) de \(P_3\) do exemplo, \(x=(t,3)\) com \(t\in[-1,1]\). De \(x\), o caminho visita \(P_1\) e \(P_2\) e volta. Os contatos ficam nos cantos \((-2,1)\) e \((2,1)\): o elo do meio mede pelo menos 4, e cada canto é o ponto da sua caixa mais próximo de \(x\). Então
\[
F(t)=\sqrt{(t+2)^2+4}+4+\sqrt{(2-t)^2+4},\qquad
f'(t)=\frac{t+2}{\sqrt{(t+2)^2+4}}-\frac{2-t}{\sqrt{(2-t)^2+4}},
\]
que é exatamente \eqref{eq:derivada} com \(u_{\text{in}}=(t-2,2)/\norm{\cdot}\) e \(u_{\text{out}}=(-2-t,-2)/\norm{\cdot}\). A derivada vale \(\mp0{,}3849\) nas pontas e zero em \(t=0\), onde \(F=4+4\sqrt2\) (\cref{fig:ancora}).

\begin{figure}[ht]\centering\figancora
\caption{A função âncora. (a) O caminho ancorado em \(x=(0{,}6;3)\): sai de \(x\), visita \(P_1\) e \(P_2\) e volta, com comprimento \(F(0{,}6)\approx9{,}7216\); tracejado, o ciclo ótimo. (b) \(F\) é convexa ao longo da aresta, com derivada negativa na ponta esquerda e positiva na direita; o mínimo, em \(t=0\), é \(\OPT\).}\label{fig:ancora}
\end{figure}

\paragraph{A busca exata} Em cada trecho, o fluxo exato testa primeiro as duas pontas: cada ponta é testada contra o problema global e contra o problema restrito em que a região da âncora é só o trecho fechado; o certificado completo trata as pontas não suaves. Se nenhuma ponta resolve, a convexidade garante derivada negativa à esquerda e positiva à direita, e uma bissecção no sinal de \eqref{eq:derivada} encontra o ponto estacionário. Para terminar em aritmética exata, o código reconstrói, por frações contínuas, o racional mais simples do intervalo: num tipo combinatório fixo, o comprimento desdobrado é uma soma de distâncias a imagens racionais, ou \(\norm{(A-I)x+c}\) com \(A,c\) racionais, cujo quadrado é quadrático no parâmetro da aresta; o minimizador é racional, com algum denominador \(Q\). Um intervalo de largura menor que \(1/Q^2\) contém um único racional de denominador até \(Q\), e então a bissecção seguida da reconstrução acha o minimizador em \(O(1+\log Q)\) passos. A largura do intervalo nunca é um critério de aceitação: só o certificado global aceita.

\begin{algorithm}[ht]\small
\caption{Redução à fronteira (fluxo exato; esquema).}\label{alg:fronteira}
\begin{algorithmic}[1]
\Function{CicloExato}{$P_1,\dots,P_k$}
  \IfThen{$\bigcap_i P_i\ne\emptyset$}{\Return o ciclo zero num ponto comum}
  \State tente as propostas da construção rápida (\cref{sec:rapida}); se uma passa no certificado, \Return
  \For{cada trecho $[a,a+e]$ das fronteiras subdivididas}
    \For{cada ponta $x$}
      \State resolva o caminho ancorado ($s=t=x$, mapas de Dror); \IfThen{certificado global ou restrito passa}{use $x$}
    \EndFor
    \State bissecção no sinal de $(u_{\text{in}}-u_{\text{out}})\cdot e$; reconstrução do racional mais simples
    \IfThen{o certificado global aceita o ciclo}{\Return ele}
  \EndFor
\EndFunction
\end{algorithmic}
\end{algorithm}

\subsection{O caminho rápido: contatos ativos}\label{sec:rapida}
A redução à fronteira é o último recurso, não o caminho usual. Quase toda chamada fecha antes, com uma \emph{proposta de tipo combinatório}: para cada contato, se ele está num vértice, no interior de uma aresta ou ``atravessando'' a região (em linha reta entre os vizinhos).

\begin{enumerate}[leftmargin=*,itemsep=2pt]
\item \textbf{Varreduras de coordenada.} Partindo dos centroides (ou dos contatos herdados do pai, \cref{sec:relaxacao}), cada contato minimiza os seus dois elos com os vizinhos fixos: atravessa se o segmento entre os vizinhos cruza a região; senão, testa os vértices e as reflexões nas arestas. Contatos consecutivos coincidentes podem se mover juntos sobre a interseção das suas regiões (um ponto, um segmento ou um polígono), o que tira blocos de arestas compartilhadas de armadilhas como a da \cref{sec:janelas}.
\item \textbf{Solução fechada do tipo.} Com contatos fixos em vértices (``pinos''), as reflexões entre dois pinos consecutivos se desdobram num segmento reto, cuja interseção com cada aresta dá os contatos. No exemplo, o contato de \(P_3\) fica entre os pinos \((2,1)\) e \((-2,1)\); refletindo \((-2,1)\) na reta \(y=3\) obtemos \((-2,5)\), e o segmento de \((2,1)\) a \((-2,5)\) corta \(y=3\) em \((0,3)\). Sem pinos, ancora-se numa aresta \(x=a+te\): o desdobramento dá um comprimento \(\norm{(A-I)x+c}\), cujo quadrado é quadrático em \(t\), e o ponto estacionário é racional.
\item \textbf{Pivôs.} Se o desdobramento põe um contato fora da sua aresta, ele é preso à ponta que bloqueia e a cadeia é refeita; cada pivô acrescenta um pino, então há no máximo \(k\) por varredura. O tipo resultante alimenta a varredura seguinte, até \(k+1\) varreduras ou um tipo repetido.
\end{enumerate}
Toda proposta passa pelo certificado; um tipo errado nunca autoriza \code{Optimal}. O núcleo é o mesmo nas duas aritméticas (\code{CycleRefinement<Scalar>}). Um tipo vindo do pai (opção \code{features}, usada nas campanhas) é reparado só nos contatos que mudaram, antes da varredura completa.

\subsection{O estágio em \texttt{binary64}}\label{sec:binary64}
Desde 08/10, uma chamada com gap pedido \(\tau>0\) ou corte finito tenta primeiro um estágio que prova \(L\le\OPT\le U\) com arredondamento dirigido em \texttt{binary64}, sem aritmética racional. É a versão de ciclo do oráculo de caminho, e reaproveita as suas peças~\cite{oraculo}.

\paragraph{Normalização} Cada região é normalizada em \texttt{binary64} para o mesmo conjunto que o fluxo exato aceitaria: duplicatas removidas, um vértice é um ponto, dois são um segmento; nos polígonos, cada sinal de curva é decidido exatamente (intervalo que exclui zero, senão o determinante inteiro de 128 bits), a fronteira precisa girar uma vez e vértices colineares são removidos. Um sinal indeciso devolve \code{Unsupported}, e a chamada vai para o fluxo exato.

\paragraph{Propostas} Um ponto comum de todas as regiões (o ciclo zero), cada candidato do núcleo de construção em \texttt{double} e, se nenhum fecha, o polimento abaixo. Nenhuma proposta é confiada.

\paragraph{Limite superior} Cada contato precisa estar provadamente na sua região: num ponto, é o próprio ponto; num segmento \([a,b]\), um contato que não está exatamente nele vira o ponto exato \(a+t(b-a)\), com \(t\in[0,1]\) em \texttt{binary64}, carregado como uma caixa intervalar; num polígono, cada sinal de aresta é decidido por um filtro estático, por intervalo ou pelo determinante inteiro. Um contato não provado é puxado para a média dos vértices (\(2^{-45},2^{-40},\dots,1\)) e testado de novo. \(U\) é o extremo superior da soma intervalar dos elos.

\paragraph{Limite inferior} \(L=\max(0,D(u))\) para cada \(u\) proposto, avaliado de modo que o valor calculado nunca passe do exato. Cada vetor é arredondado para \texttt{binary64} e tem norma provada \(\le1\). Cada suporte é o mínimo em \texttt{binary64} alargado por um limite \emph{a priori} do seu erro de arredondamento (\cref{lem:apriori}, na \cref{sec:insercao}), e as somas são arredondadas para baixo. Os elos curtos, de comprimento até \(\max(32\epsilon\cdot\text{escala},\tau/(16k))\), recebem quatro propostas: a própria direção, a do elo longo mais próximo de cada lado (ciclicamente) ou zero. O comprimento serve só para escolher propostas; nunca é tolerância.

\paragraph{Polimento} Se nenhum candidato fecha o gap, um método de Newton com barreira logarítmica minimiza o comprimento suavizado
\[
f_\mu(y)=\sum_{i=1}^k\sqrt{\norm{x_{i+1}-x_i}^2+\mu^2}\;-\;\mu\sum_i\sum_{f\in F_i}\log(n_f\cdot y_i-o_f),
\]
em coordenadas escaladas, com uma barreira por aresta (num segmento \(x=a+te\), a barreira é \(-\mu(\log t+\log(1-t))\); um ponto é fixo). Num ponto estacionário, com \(u_i=d_i/\sqrt{\norm{d_i}^2+\mu^2}\), vale
\begin{equation}\label{eq:perdabarreira}
D(u)\;\ge\;\sum_i\norm{d_i}-\mu\,(k+F),\qquad F=\textstyle\sum_i\abs{F_i},
\end{equation}
a mesma proposição do caminho, com \(k\) elos no lugar de \(m+1\). Ela só justifica a escolha de \(\mu\) final (metade do gap pedido); os limites devolvidos são sempre os da prova acima. O sistema de Newton é tridiagonal por blocos \(2\times2\), como no caminho, exceto pelo elo de fechamento, que acopla a primeira e a última variável. O solver elimina os \(k-1\) primeiros blocos para o lado direito e para as duas colunas do último bloco, resolve o complemento de Schur \(2\times2\) da última variável e faz a substituição de volta: \(O(k+F)\) operações por iteração, sem pivoteamento (a barreira mantém a hessiana definida positiva).

\begin{teorema}[Correção dos limites do estágio]\label{teo:binary64}
Se o ambiente é IEEE \texttt{binary64} com arredondamento ao mais próximo (o código verifica), os limites \((L,U)\) devolvidos pelo estágio satisfazem \(L\le\OPT\le U\), qualquer que seja a proposta. Com \(U\) finito, ele limita o comprimento de um ciclo de pontos exatos \(z_i\in P_i\).
\end{teorema}
\begin{proof}[Ideia da prova]
A normalização aceita exatamente as regiões que o fluxo exato aceita e não muda os conjuntos. Cada caixa de contato contém um ponto exato da sua região, então os pontos formam um ciclo viável, e a soma intervalar limita o seu comprimento por cima. Cada \(u_i\) tem norma provada \(\le1\), e cada termo de \eqref{eq:dualciclo} é limitado por baixo pelo mínimo calculado menos o erro \emph{a priori}; a soma arredondada para baixo é \(\le D(u)\le\OPT\). Nenhum passo usa a otimalidade da proposta. A prova completa está em~\cite{ciclo}.
\end{proof}

\subsection{O oráculo completo e o seu contrato}\label{sec:oraculociclo}
O \cref{alg:oraculo} junta as peças. A busca chama \code{tpp\_convex\_solve\_cycle\_double} com o gap \(\tau\) e o corte \(c\) da mesma regra do caminho~\cite{oraculo}: \(\tau=\max(g/4,\rho\,U^\star)\) nas chamadas comuns e \(g/4\) nos refinamentos, com \(g\) o gap do nó e \(\rho=10^{-6}\), e \(c=U^\star-g\). Um estágio fechado devolve \code{GapClosed} (ou \code{CertifiedBound}, se \(L\ge c\)) sem teste KKT e sem nenhuma etapa racional. Um estágio aberto passa ao fluxo exato, e o resultado fica com o maior dos dois limites inferiores (ambos limitam o mesmo ótimo).

\begin{algorithm}[ht]\small
\caption{O oráculo de ciclo como a busca o chama (esquema).}\label{alg:oraculo}
\begin{algorithmic}[1]
\Function{Ciclo}{$P_1,\dots,P_k$; contatos e tipo herdados; $\tau$; $c$}
  \If{$\tau>0$ ou $c<\infty$}\Comment{estágio \texttt{binary64}}
    \State normalize; proponha (ponto comum, construção, depois polimento)
    \For{cada proposta}
      \State prove $U$ (contatos nas regiões) e $L$ (\eqref{eq:dualciclo} com erro \emph{a priori})
      \IfThen{$U-L\le\tau$ ou $L\ge c$}{\Return $(L,U)$, caminho}
    \EndFor
  \EndIf
  \State fluxo exato: construção racional, certificado KKT, redução à fronteira (\cref{alg:fronteira})
  \State \Return o maior $L$ e o caminho de menor $U$
\EndFunction
\end{algorithmic}
\end{algorithm}

Há ainda um passo no adaptador da busca (\code{solve\_relaxation}): se o intervalo devolvido continua mais largo que \(\tau\) e abaixo do corte, o ciclo é resolvido de novo em aritmética racional, partindo dos contatos já verificados. Isso importa em pontos e segmentos, que não têm interior para onde arredondar: contatos arredondados podem ficar bem acima do ótimo da relaxação, e antes a busca guardava nós cujo gap nunca fechava. Na coleção da Paula, chamadas que chegam a qualquer etapa racional são raras: 7 em 99,1 milhões (3 que o polimento deixa abertas e 4 interrompidas pelo prazo).

\begin{cuidado}[O que depende do quê]
Como no caminho, os limites não dependem da correção da construção: nem dos mapas de último passo com interseção (cuja prova geral continua em aberto~\cite{tutorial}), nem do núcleo de contatos ativos, nem do polimento. Desses depende que o oráculo \emph{feche} o gap pedido, e quanto tempo isso leva; um erro neles alarga o intervalo, mas nunca o torna falso.
\end{cuidado}

\subsection{Pontos e segmentos}
As instâncias da Paula têm pontos (clusters de um elemento) e segmentos. Eles entram no ciclo como regiões convexas sem interior, com três cuidados. Na interseção comum, um segmento é limitado pela sua reta \emph{e} pelos seus extremos, e um ponto é testado por pertencimento: os semiplanos das ``arestas'' de um segmento descrevem uma reta infinita, e os de um ponto, o plano inteiro. Os contatos num segmento são carregados como \(a+t(b-a)\) com \(t\) em \texttt{binary64}, o que os mantém no segmento exato. E a construção devolve o próprio ponto como contato de uma região-ponto: uma versão anterior tratava o ponto como o plano inteiro e propunha contatos ao lado dele, que o certificado sempre rejeitava.

%==========================================================================
\section{O branch-and-bound}\label{sec:bnb}

\subsection{Uma busca, duas topologias}
\code{tpp\_nonconvex\_tspn\_solve} chama a mesma função de busca do TPP de ordem livre (\code{solve\_normalized\_unordered\_tpp}) com uma flag de ciclo. Um nó guarda uma sequência cíclica parcial; cada elemento é o fecho convexo de uma região original ou uma peça da sua decomposição. O limite do nó é o ótimo da relaxação (o ciclo de ordem fixa por esses convexos) ou, antes de resolvê-la, o limite herdado e o da triagem. Fila de melhor limite, mergulhos, decomposição sob demanda, escolha da região mais distante, verificação de visitas, normalização afim e contabilidade da fronteira são os mesmos. Esta seção descreve só o que muda; o contrato completo está em~\cite{tspnmd}.

\subsection{Quando há um ponto, o TSPN vira TPP}\label{sec:ancora}
\begin{proposicao}[Âncora de ponto]\label{prop:ancoraponto}
Se \(R_a=\{p\}\) é um ponto, o TSPN de \(R_1,\dots,R_n\) tem o mesmo valor que o TPP de ordem livre de \(p\) a \(p\) pelas outras regiões, e os ótimos se correspondem.
\end{proposicao}
\begin{proof}
Todo tour passa por \(p\). Rotacionado para começar e terminar em \(p\), ele é um caminho de \(p\) a \(p\) que visita as outras regiões, com o mesmo comprimento. Reciprocamente, todo caminho de \(p\) a \(p\) que visita as outras regiões é um tour fechado que visita todas, inclusive \(R_a\).
\end{proof}
Então, quando há um ponto (66 das 235 instâncias da Paula, entre elas 6 das 15 ainda abertas), o solver usa o TPP de ordem livre com \(s=t=p\), inteiro: o oráculo de caminho com extremos fixos, as \(m+1\) posições de inserção e o limite de inserção do caminho. Um tour inicial fornecido é rotacionado para começar no ponto. Na campanha de 06/10, a âncora fechou 45 instâncias com ponto contra 42 sem ela, e foi 2,67 vezes mais rápida. O resto desta seção trata do caso sem pontos.

\subsection{A raiz}\label{sec:raiz}
Por padrão, a raiz é o fecho da região 0, com limite zero. Um ciclo por uma só região é qualquer ponto dela; fixar a região 0 como ``primeira'' elimina a simetria de rotação sem impor um ponto artificial, porque todo tour visita a região 0 em algum lugar e pode ser lido a partir dessa visita.

Com a opção \code{root} (usada nas campanhas), a raiz é uma tripla separada: o par de regiões mais distantes entre si, \(a\) e \(b\), e a região \(c\) que maximiza \(\dist(a,c)+\dist(b,c)\). Começar com três regiões não perde nenhum tour: há uma só ordem cíclica não orientada de três rótulos, e todas as inserções seguintes continuam representadas. As distâncias só ordenam; não são limites. A vantagem é começar com um ciclo que já tem comprimento, em vez de limite zero.

\subsection{Ramificação}\label{sec:ramificacao}
Como no caminho, a busca escolhe a região original não visitada mais distante do tour relaxado. Se ela não está na sequência, o seu fecho é inserido em cada posição cíclica, inclusive no elo de fechamento: \(m\) posições, contra \(m+1\) no caminho (\cref{fig:lacunas}). Com \(m=2\), os dois filhos são reversões do mesmo triângulo, e só um é gerado; a reversão preserva comprimento e visitas. Se o fecho já está na sequência, ele é substituído por cada peça da decomposição convexa.

\begin{figure}[ht]\centering\figlacunas
\caption{Posições de inserção (esquema). (a) No caminho, os extremos são distintos e há \(m+1\) lacunas. (b) No ciclo, o elo de fechamento é uma lacuna como as outras, e há \(m\).}\label{fig:lacunas}
\end{figure}

\paragraph{Por que a ramificação é completa} Fixe um tour viável e uma visita a cada região já representada; a visita à região 0 serve de origem. As visitas às regiões representadas ocorrem numa ordem cíclica, e a visita à região nova cai em alguma das lacunas dessa ordem, inclusive a de fechamento. Uma visita a uma região não convexa pertence a alguma peça. Omitir regiões e trocar regiões pelos fechos relaxa o problema, então o ótimo de cada relaxação limita por baixo todas as completações do nó. Como no caminho, cruzamentos do tour parcial e visitas repetidas a uma região não autorizam poda; os contraexemplos registrados continuam valendo~\cite{contraexemplos}.

\subsection{A relaxação e os incumbentes}\label{sec:relaxacao}
Cada filho recebe como ponto de partida os contatos do pai, com o contato da região nova dado por \code{best\_contact} entre os dois vizinhos (\cref{sec:melhorcontato}) e, com \code{features}, o tipo combinatório do pai com a posição nova marcada como desconhecida. Isso é só proposta; o certificado valida tudo. O oráculo também recebe o corte do incumbente: um filho cujo limite certificado alcança o corte volta com \code{CertifiedBound} e é podado sem que se construa o seu ótimo. Um tour relaxado que visita todas as regiões originais (verificação de visitas com tolerância \(10^{-8}\)) é um incumbente global.

\subsection{A triagem por limites de inserção}\label{sec:insercao}
Esta é a peça que mais pesou no desempenho recente. A ideia: o dual do pai já é quase um dual do filho. Inserir uma região numa lacuna muda só três termos de \eqref{eq:dualciclo}, então o limite do filho sai de algumas operações sobre os vértices de três regiões, sem chamar o oráculo.

\paragraph{A conta} Seja \(q_1,\dots,q_m\) o tour relaxado do pai, \(u_1,\dots,u_m\) as direções dos seus elos (\(u_j\) vai de \(q_j\) a \(q_{j+1}\); um elo de comprimento zero recebe a direção de um vizinho) e \(D_{\text{pai}}=\sum_i\supp_i(u_{i-1}-u_i)\), com \(\supp_i(c)=\min_{v\in P_i}(v-r)\cdot c\). Para inserir a região \(R\) na lacuna \(j\), entre \(q_j\) e \(q_{j+1}\), escolhemos um contato proposto \(c=\code{best\_contact}(q_j,q_{j+1},R)\) e as direções \(a\) de \(q_j\) a \(c\) e \(b\) de \(c\) a \(q_{j+1}\). O filho tem \(m+1\) elos; usamos para ele as direções do pai, trocando \(u_j\) pelo par \((a,b)\). Então
\begin{equation}\label{eq:insercao}
D_{\text{filho}}=D_{\text{pai}}
\;-\;\supp_j(u_{j-1}-u_j)-\supp_{j+1}(u_j-u_{j+1})
\;+\;\supp_j(u_{j-1}-a)+\supp_R(a-b)+\supp_{j+1}(b-u_{j+1}).
\end{equation}
Por dualidade fraca, \(D_{\text{filho}}\) limita por baixo o ótimo de \emph{qualquer} completação do filho, por pior que seja o contato proposto: a escolha de \(c\) só afeta a força do limite. O limite do filho é \(\max(L_{\text{pai}},D_{\text{filho}})\), e um filho com limite \(\ge U^\star-g\) é descartado antes do oráculo. Os suportes que não mudam são somados uma vez para todos os irmãos.

\paragraph{Exemplo} Inserimos \(R=[-0{,}5;0{,}5]\times[-3,-2]\) na lacuna entre \(P_1\) e \(P_2\) do exemplo das três caixas (\cref{fig:insercao}). Os contatos do pai são \(q_1=(-2,1)\) e \(q_2=(2,1)\), e o contato proposto é \(c=(0,-2)\), o ponto do topo de \(R\) que minimiza \(\norm{q_1-c}+\norm{c-q_2}\). Então \(a=(2,-3)/\sqrt{13}\) e \(b=(2,3)/\sqrt{13}\), e os termos de \eqref{eq:insercao} estão na \cref{tab:insercao}.
\begin{table}[ht]\centering\small
\begin{tabular}{@{}lccc@{}}\toprule
termo & coeficiente & mínimo em & valor\\\midrule
\(\supp_1(u_3-a)\) & \((-1{,}2618;\;0{,}1249)\) & \((-2,-1)\) & \(2{,}3987\)\\
\(\supp_R(a-b)\) & \((0;\;-1{,}6641)\) & aresta \(y=-2\) & \(3{,}3282\)\\
\(\supp_2(b-u_2)\) & \((1{,}2618;\;0{,}1249)\) & \((2,-1)\) & \(2{,}3987\)\\
\(\supp_3(u_2-u_3)\) (inalterado) & \((0;\;1{,}4142)\) & aresta \(y=3\) & \(4{,}2426\)\\\bottomrule
\end{tabular}
\caption{Os termos de \eqref{eq:insercao} no exemplo, com \(r=0\): três suportes mudam e um é herdado do pai.}\label{tab:insercao}
\end{table}

A soma dá \(D_{\text{filho}}\approx12{,}3682\), contra \(D_{\text{pai}}=4+4\sqrt2\approx9{,}6569\); o ganho, \(2{,}7113\), mede o desvio necessário para passar por \(R\) a partir dessa lacuna. O ótimo do filho, conferido pelo solver, é \(2\sqrt{41}\approx12{,}8062\): o limite de inserção fica 3,4\% abaixo dele, sem nenhuma chamada ao oráculo. Se o incumbente fosse menor que 12,37, o filho seria podado ali mesmo.

\begin{figure}[ht]\centering\figinsercao
\caption{O limite de inserção. (a) O ciclo do pai (tracejado), a caixa \(R\) inserida na lacuna entre \(P_1\) e \(P_2\), o contato proposto \(c\) e as direções \(a,b\) que substituem \(u_1\). O limite resultante é \(12{,}3682\). (b) Um ótimo do filho, de comprimento \(2\sqrt{41}\approx12{,}8062\). Este é o ótimo simétrico; o solver devolveu outro do mesmo comprimento, da mesma família de paralelogramos.}\label{fig:insercao}
\end{figure}

\paragraph{A prova em \texttt{binary64}} Até 09/10, \eqref{eq:insercao} era avaliada em aritmética racional, com direções racionais de norma provada; depois que o oráculo passou a \texttt{binary64}, essa triagem dominava o tempo. Hoje ela é provada em \texttt{binary64} com três ingredientes:
\begin{itemize}[leftmargin=*,itemsep=2pt]
\item \textbf{Direções.} Cada vetor é a direção calculada em \texttt{binary64}, encurtada por \(2^{-49}\). Um quadrado da norma calculado até \(1-2^{-50}\) prova que a norma exata é no máximo 1 (\cref{lem:apriori}); se não, quadrados intervalares decidem; um vetor não provado vira zero, que é sempre viável. Um elo novo de comprimento zero pode ficar nulo ou herdar \(u_j\); vale o que der o ganho maior.
\item \textbf{Suportes.} Cada suporte é o mínimo calculado alargado pelo erro \emph{a priori} do lema abaixo; uma normal exatamente nula dá suporte exatamente zero.
\item \textbf{Somas.} Arredondadas para fora, um valor representável por operação; o ganho é o extremo inferior do seu intervalo.
\end{itemize}

\begin{lema}[Erro \emph{a priori} de um termo de suporte]\label{lem:apriori}
No modelo padrão de arredondamento (erro relativo até \(\mathbf u=2^{-53}\) por operação, com \emph{underflow} gradual), com \(v\), \(o\) e \(n\) em \texttt{binary64} e \(d=v-o\) calculado em \texttt{binary64}, o valor calculado de \(t=d_x n_x+d_y n_y\) está a no máximo \(4{,}0000001\,\mathbf u\,(\abs{d_x}\abs{n_x}+\abs{d_y}\abs{n_y})+2^{-1074}\) do valor exato de \((v-o)\cdot n\), inclusive se o compilador usar FMA. Então o mínimo calculado sobre os vértices, menos \(E=2^{-50}(m_x\abs{n_x}+m_y\abs{n_y})+2^{-1070}\), com \(m_x,m_y\) as maiores diferenças calculadas, é no máximo o suporte exato.
\end{lema}
O lema troca intervalos dirigidos por uma única correção por suporte, três a quatro vezes mais barata, e a prova está em~\cite{ordemlivre}. O resultado é um limite de dualidade fraca para o filho, como o racional, com vetores ligeiramente diferentes; não é uma tolerância. Na regressão SoCG, a busca ficou idêntica (as mesmas chamadas nos 48 casos) e a triagem caiu de 22,4 para 0,15 s.

\begin{ideia}[Limite de inserções múltiplas (opcional)]
A triagem acima olha uma região de cada vez. A opção \code{-{}-multi-insertion-bound} eleva o limite de um nó com os ganhos de \emph{todas} as regiões ausentes: lacunas não adjacentes mudam termos disjuntos e os ganhos se somam; duas lacunas adjacentes perdem no máximo a largura da região que compartilham na direção do pai. Um dual de cobertura atribui preços às regiões, de modo que, qualquer que seja a lacuna em que cada uma acabe, a soma dos preços é paga; os preços gulosos são conferidos com somas arredondadas para cima antes de valer. É provado em \texttt{binary64} como a triagem, mas custa \(O(m\,k)\) contatos e suportes por nó; desde que o oráculo ficou barato, ele não compensa no tempo e está desligado. Com a memória como limite (\cref{sec:memoria}), vale reavaliá-lo.
\end{ideia}

\subsection{A política de busca}\label{sec:politica}
\paragraph{Padrão: melhor limite com mergulhos} A fila expande o nó de menor limite, e a cada expansão mergulha pelo melhor filho até uma folha (\code{dive\_interval} 1), para achar incumbentes cedo. Os filhos que passam pela triagem são resolvidos logo (avaliação imediata), em ordem; com \code{-{}-threads N}, até \(N\) irmãos em paralelo. É a política que fecha as instâncias fáceis: 209 das 235 da Paula em 60 s, sem tour inicial.

\paragraph{Preguiçosa, para o limite inferior} Com \code{-{}-cycle-optimization lazy} (ciclos) e \code{-{}-lazy-oracles} (caminhos ancorados num ponto), os filhos entram na fila com o limite da triagem, e o oráculo só é chamado quando um deles é escolhido para expansão; um incumbente melhor pode podá-lo antes. Isso economiza chamadas, mas atrasa os incumbentes. Com um incumbente bom vindo da ILS, os mergulhos também deixam de pagar, e a combinação \code{-{}-primal-ils 0.5 -{}-cycle-optimization lazy -{}-lazy-oracles -{}-dive-interval 0} é a que mais sobe o limite inferior: nas 38 instâncias abertas aos 60 s, 600 s dão gap mediano de 2,7\% e fecham 17. Sem incumbente, a política padrão continua melhor: em 60 s, a preguiçosa sem mergulhos também fecha as 209, mas 0,79 vez a velocidade da padrão em média geométrica, e só tirar os mergulhos perde 6.

\paragraph{Portfólio} \code{-{}-portfolio} roda duas buscas completas em paralelo: a de melhor limite e uma DFS/BFS inspirada no solver de Fekete et al., que desce pelo filho mais barato e volta ao melhor nó da fronteira depois de uma poda. Elas trocam incumbentes validados, e o maior dos dois limites inferiores é global. É opcional.

\subsection{A memória da fronteira}\label{sec:memoria}
Com o oráculo em \texttt{binary64}, a fronteira de melhor limite cresce na proporção da vazão: em 100pr1002, 60 s chegam a 4,2 milhões de nós e 3–3,8 GB. Desde 09/10 há um orçamento (\code{-{}-max-frontier-mib}, padrão 4 GiB) com dois níveis. Ao passar do limite, todos os nós menos o quarto de menor limite perdem caminho, contatos e tipo, que são recalculados se o nó sair da fila (``despir''). Se isso não basta, a metade de maior limite é descartada, e o menor limite descartado, \(D\), vira um piso do limite inferior global: \(L=\min(\text{fronteira},D)\). Isso é válido, porque toda completação de um nó descartado custa pelo menos \(D\).

O piso tem uma consequência que só medimos em 10/10 (\cref{fig:piso}): quando o mínimo da fronteira alcança \(D\), o limite inferior \emph{congela}, e o resto da execução não pode mais subi-lo. Com 256 MiB, 100i1000-410 congela aos 2 min em 4,88\%; com 6 GiB, no mesmo tempo, continua caindo até 3,64\% aos 5 min. Nessas condições, um oráculo mais rápido só enche a fronteira mais cedo e congela no mesmo valor; o que ajuda é mais memória por execução, menos memória por nó, ou não perder os nós descartados (\cref{sec:limites}).

Na armijo (128 GB), com 22 GiB por execução e 4 h, seis das oito instâncias abertas com gap até 5\% fecharam, em 22 a 78 min, com picos de fronteira de 3,9 a 14,7 GiB; as rodadas anteriores, com 1,5 GiB, tinham congelado em 2,4–4,3\%. Isso levou a três mudanças de 12/10, que não alteram a busca. Os irmãos compartilham uma única cópia dos contatos e do tipo combinatório do pai, em vez de cada um copiá-los como ponto de partida, e cada filho guarda só o seu contato novo; o caminho, as features e o dual do nó passam a custar um ponteiro enquanto vazios. Um nó de ciclo na fila de \code{lazy} cai de \(\approx\)900 para \(\approx\)285 B, e um de caminho ancorado, de 253 para 197 B. E a busca para com \code{memory\_limit} assim que o limite congela no piso, em vez de rodar até o prazo.

\begin{figure}[ht]\centering\figpiso
\caption{O piso da fronteira. O mesmo B\&B (preguiçoso, sem mergulhos, a partir do mesmo tour) em 100i1000-410, com dois orçamentos. Com 256 MiB, o limite inferior para de subir aos 120 s; com 6 GiB (pico de 5 GiB), continua. Dados em \code{dados.tex}.}\label{fig:piso}
\end{figure}

\subsection{O contrato numérico}
A busca trabalha em coordenadas normalizadas (centro e divisor), e os limites voltam às coordenadas de entrada com o erro da normalização descontado. Um resultado \code{exact} significa que o gap pedido fechou, \(U-L\le g_a+g_r\abs U\), com limites provados das relaxações e verificação de visitas com tolerância \(10^{-8}\); não é uma solução racional sem erro do TSPN não convexo completo. As campanhas usam \(g_a=0\) e \(g_r=10^{-6}\), e cada tour devolvido é validado de forma independente, com testes de interseção e pertencimento em racionais binários (\code{validate\_cycle}, tolerância \(10^{-7}\)). Limites de tempo, de chamadas e de memória devolvem o incumbente e os limites com um status incompleto (\code{time\_limit}, \code{call\_limit}, \code{memory\_limit}); nenhum deles é chamado de ótimo.

%==========================================================================
\section{A busca local iterada}\label{sec:ils}

\subsection{Para que serve}
A ILS (\code{-{}-primal-ils F}) gasta a fração \(F\) do tempo restante procurando um bom tour antes do \emph{branch-and-bound}. Ela só mexe no limite superior: todo tour que ela encontra passa pela mesma verificação de visitas e de comprimento, e o limite inferior e o certificado continuam os do B\&B. Serve em dois casos: nas instâncias grandes, onde o B\&B não fecha e o que se entrega é um bom tour com gap provado; e na política preguiçosa, que precisa de um incumbente bom desde o início. No TPP de ordem livre com extremos fixos, ela não acelerou o B\&B e fica desligada.

\subsection{Representação}
Um tour é uma ordem das \(n\) regiões originais e um contato por região, \emph{dentro da própria região original}, que pode ser não convexa. Todo movimento mantém essa invariante, então todo tour é viável por construção; mesmo assim, só um tour validado vira incumbente.

\subsection{O subproblema de um contato}\label{sec:melhorcontato}
Quase todo movimento termina reotimizando contatos com os vizinhos fixos: dado \(\ell\) e \(r\), achar \(p\in P\) que minimiza \(\norm{\ell-p}+\norm{p-r}\). É o mesmo subproblema da descida por coordenadas (BCD) da Paula~\cite{paula}. A função \code{best\_contact} o resolve assim: se \(\ell\) ou \(r\) está em \(P\), devolve esse ponto (custo \(\norm{\ell-r}\), o mínimo possível); senão, para cada aresta, calcula o melhor ponto da aresta e fica com o melhor.

\begin{proposicao}[Uma fórmula para reflexão e cruzamento]\label{prop:melhorcontato}
Seja uma aresta \([a',a'+e]\) e sejam \(h_\ell,h_r\) as distâncias de \(\ell\) e \(r\) à reta da aresta, com \(h_\ell+h_r>0\). O ponto da reta que minimiza \(\norm{\ell-p}+\norm{p-r}\) tem como parâmetro a média das projeções de \(\ell\) e \(r\) com pesos \(h_r/(h_\ell+h_r)\) e \(h_\ell/(h_\ell+h_r)\), tanto com \(\ell\) e \(r\) do mesmo lado da reta (reflexão) quanto em lados opostos (cruzamento). Limitado a \([0,1]\), esse parâmetro dá o melhor ponto da aresta.
\end{proposicao}
\begin{proof}
Em lados opostos, o mínimo é o cruzamento do segmento \(\ell r\) com a reta. Do mesmo lado, refletindo \(r\) na reta, o mínimo é o cruzamento do segmento de \(\ell\) ao refletido. Nos dois casos, o ponto divide o segmento de \(\ell\) ao alvo (\(r\) ou o refletido) na razão \(h_\ell:h_r\), por semelhança de triângulos, e a projeção na reta é linear; daí os pesos. A função é convexa no parâmetro, então limitar o minimizador irrestrito ao intervalo dá o minimizador restrito.
\end{proof}
Com \(\ell=(0,0)\), \(r=(6,1)\) e a aresta de \((1,2)\) a \((5,2)\), \(h_\ell=2\) e \(h_r=1\), e o ponto é \((4,2)\); com \(r=(6,3)\), do outro lado, o ponto é o mesmo (\cref{fig:melhorcontato}). A rotina é exata (em aritmética exata) também para polígonos não convexos: a função é convexa no plano, e se o seu mínimo em \(P\) estivesse no interior de \(P\) seria um mínimo local, portanto global, e estaria no segmento \(\ell r\), que não toca \(P\) nesse caso. Logo o mínimo está na fronteira, e cada aresta é tratada exatamente.

\begin{figure}[ht]\centering\figmelhorcontato
\caption{O subproblema de um contato numa aresta horizontal. (a) Com \(\ell\) e \(r\) do mesmo lado, o melhor ponto é o da reflexão (tracejado: o refletido \(r'\)). (b) Em lados opostos, é o cruzamento. A mesma média ponderada das projeções, com pesos dados por \(h_\ell=2\) e \(h_r=1\), dá \(p=(4,2)\) nos dois casos.}\label{fig:melhorcontato}
\end{figure}

\subsection{Solução inicial}
A heurística inicial é a mesma do B\&B sem ILS. A partir de um vértice da região 0, escolhe sempre o vértice mais próximo de uma região ainda não visitada; aplica até 10 passadas de 2-opt; depois, até 8 passadas alternando uma varredura de \code{best\_contact} em todos os contatos e um 2-opt. A ILS parte do melhor tour assim obtido.

\subsection{A busca local}
O \cref{alg:ils} mostra o laço. A busca local repete, até não melhorar (no máximo 64 rodadas):
\begin{enumerate}[leftmargin=*,itemsep=2pt]
\item \textbf{Varredura de contatos.} Cada contato vira \code{best\_contact} entre os seus vizinhos.
\item \textbf{2-opt.} Inverte as posições \(i..j\) com os contatos fixos, se isso encurta o tour.
\item \textbf{Or-opt.} Move um bloco de 1 a 3 posições consecutivas para a melhor lacuna; os contatos das pontas do bloco são reotimizados para os novos vizinhos (um contato, se o bloco tem uma posição). Com listas de candidatos, só se tentam lacunas vizinhas de uma das 10 regiões mais próximas (distância entre fronteiras) de uma ponta do bloco; isso dá cerca de 5 vezes mais iterações por segundo em 100 regiões.
\end{enumerate}
A troca de duas regiões existe como opção (\code{-{}-primal-ils-swap}) e foi pior.

\subsection{Perturbação e aceitação}
A perturbação é uma \emph{double bridge} aleatória nas posições (\(ABCD\to ACBD\)), seguida da busca local. A aceitação é a da Paula: o tour tentado substitui o corrente se for melhor que ele ou estiver a no máximo \(\eta\) acima do melhor, em termos relativos; \(\eta\) começa em \(0{,}01\), é multiplicado por \(0{,}95\) a cada 10 iterações sem novo melhor e, abaixo de \(10^{-4}\), volta a \(0{,}01\), com o corrente voltando ao melhor (reaquecimento). A semente depende de \(n\) e de \code{-{}-primal-ils-seed}, então execuções com sementes diferentes são independentes.

\begin{algorithm}[ht]\small
\caption{A ILS (esquema; padrões entre parênteses).}\label{alg:ils}
\begin{algorithmic}[1]
\Function{ILS}{tour inicial $T$; prazo}
  \State \Call{BuscaLocal}{$T$}; \Call{Janelas}{$T$}; $T^\star\gets T$; $\eta\gets0{,}01$
  \While{houver tempo}
    \State $T'\gets$ \Call{DoubleBridge}{$T$}; \Call{BuscaLocal}{$T'$}
    \IfThen{$\abs{T'}<(1+0{,}01)\abs{T^\star}$}{\Call{Janelas}{$T'$}}\Comment{reotimização exata, janelas de 8}
    \IfThen{$\abs{T'}<\abs T$ ou $\abs{T'}\le(1+\eta)\abs{T^\star}$}{$T\gets T'$}
    \If{$\abs{T'}<\abs{T^\star}$} $T^\star\gets T'$; publique $T^\star$ como incumbente
    \Else{} a cada 10 sem melhora, $\eta\gets0{,}95\,\eta$; se $\eta<10^{-4}$: $\eta\gets0{,}01$, $T\gets T^\star$
    \EndIf
  \EndWhile
\EndFunction
\end{algorithmic}
\end{algorithm}

\subsection{A armadilha da descida por coordenadas e as janelas exatas}\label{sec:janelas}
A varredura de contatos (e o BCD da Paula) estaciona quando contatos consecutivos coincidem. O exemplo da \cref{fig:armadilha} mostra por quê. Sejam \(a=(0,0)\), \(b=(6,0)\), \(P=[1,4]\times[1,3]\) e \(Q=[2,5]\times[1,3]\), visitadas nessa ordem entre \(a\) e \(b\).

\begin{proposicao}[Todo ponto comum é um mínimo coordenada a coordenada]\label{prop:armadilha}
Se \(z\in P\cap Q\), os contatos \(p=q=z\) não melhoram por nenhuma reotimização de um contato só.
\end{proposicao}
\begin{proof}
Com \(q=z\) fixo, \(\norm{a-p}+\norm{p-z}\ge\norm{a-z}\) para todo \(p\), com igualdade em \(p=z\in P\). O mesmo vale para \(q\) com \(p=z\) fixo.
\end{proof}
Em \(z=(4,3)\), o comprimento é \(5+\sqrt{13}\approx8{,}606\). O ótimo passa por \((3,1)\in P\cap Q\) e mede \(2\sqrt{10}\approx6{,}325\), 36\% menos: todo caminho de \(a\) a \(b\) que chega a \(y\ge1\) mede pelo menos isso (reflita \(b\) na reta \(y=1\)). O B\&B de ordem livre, com extremos \(a\) e \(b\), acha esse ótimo. O rascunho da Paula prova que o BCD chega ao mínimo global sob duas hipóteses, regiões convexas e regiões consecutivas que não se tocam; as instâncias dela têm polígonos não convexos e muito sobrepostos.

\begin{figure}[ht]\centering\figarmadilha
\caption{A armadilha. Com os dois contatos em \(z=(4,3)\in P\cap Q\) (tracejado, \(5+\sqrt{13}\)), nenhuma reotimização de um contato só melhora o caminho. O ótimo (verde, \(2\sqrt{10}\)) move os dois contatos juntos até \((3,1)\).}\label{fig:armadilha}
\end{figure}

\paragraph{As janelas} Por isso a ILS tem uma etapa que a BCD não tem. Todo tour a até 1\% do melhor tem janelas de 8 posições consecutivas, a meia janela uma da outra, resolvidas \emph{exatamente} (ordem e contatos) pelo próprio B\&B de ordem livre, com os contatos fixos vizinhos como extremos \(s\) e \(t\) e as regiões originais da janela (\cref{fig:janela}). Cada janela usa gap \(10^{-9}\), até 20 mil chamadas e 0,5 s, e parte da janela atual como caminho inicial. Com regiões sobrepostas, o caminho novo pode visitar várias regiões num mesmo ponto ou ao longo de um segmento; cada região recebe o primeiro toque ao longo do caminho, na ordem desses toques, o que preserva o comprimento. A janela só é aceita se encurtar o tour. Nas 19 instâncias de validação (60 s, uma execução), a razão média geométrica entre o nosso tour e o melhor dela foi 1,0020 sem esta etapa, 1,0007 com um polimento só de contatos e 1,00015 com as janelas.

\begin{figure}[ht]\centering\figjanela
\caption{Reotimização por janela (esquema). As 8 posições da janela são resolvidas pelo B\&B de ordem livre entre os dois contatos fixos; o resto do tour não muda.}\label{fig:janela}
\end{figure}

\subsection{Relação com o ILS-BCD da Paula}
A \cref{tab:paula} compara as duas buscas componente a componente. A nossa reproduz o núcleo da dela (o subproblema de um contato, as listas de candidatos, a \emph{double bridge} e a aceitação com reaquecimento), tem menos vizinhanças e acrescenta as janelas exatas, que são o que fecha a diferença nas instâncias com regiões muito sobrepostas.

\begin{table}[ht]\centering\small
\begin{tabular}{@{}>{\raggedright\arraybackslash}p{33mm}>{\raggedright\arraybackslash}p{55mm}>{\raggedright\arraybackslash}p{56mm}@{}}\toprule
Componente & ILS-BCD (Paula) & a nossa ILS\\\midrule
Contatos & BCD: um contato por vez, subproblema exato & o mesmo subproblema (\code{best\_contact}), mais janelas exatas\\
Vizinhanças & VND com 2-opt, 3-opt, troca, inserção, realocação de blocos com inversão, Balas--Simonetti e Gutin & 2-opt, Or-opt de 1 a 3 com contatos das pontas reotimizados; troca opcional\\
Listas de candidatos & sim & sim (10 regiões por distância entre fronteiras)\\
Perturbação & \emph{double bridge} (e outra, se estagnar) & uma \emph{double bridge}\\
Aceitação & limiar \(\eta\) relativo ao melhor, com reaquecimento & a mesma\\
Parada & 3.000 iterações sem melhora ou 1.200 s de CPU & fração do tempo ou \code{-{}-primal-ils-stagnation}\\
Saída & o melhor tour & o melhor tour, e o B\&B certifica ou limita o gap\\\bottomrule
\end{tabular}
\caption{O que a nossa ILS reproduz da dela e o que difere. A nossa não tem 3-opt completo, Balas--Simonetti nem Gutin; tem as janelas exatas.}\label{tab:paula}
\end{table}

\subsection{Como a ILS e o B\&B se combinam}
A ILS roda antes da busca e publica cada novo melhor como incumbente; o B\&B começa com esse corte. Num tour inicial fornecido (\code{-{}-initial-path}), o solver valida a visita a todas as regiões e usa só o comprimento. A campanha \code{paula-campaign} separa as duas fases: 10 execuções da ILS por instância, com sementes distintas, e depois o B\&B a partir do melhor tour, com o orçamento de memória concentrado em poucas execuções (\cref{sec:memoria}).

%==========================================================================
\section{Medições}\label{sec:medicoes}
Os resumos, protocolos e limitações estão em \path{benchmarks/results-saved/README.md}~\cite{resultados}; aqui ficam os números que situam as peças. Todas as rodadas são TSPN de ciclo livre com gap relativo \(10^{-6}\), uma thread por instância e \code{cache,features,root,interval}, num Mac (Apple M4 Pro), com tours validados de forma independente. ``Fechou'' é gap numérico fechado com limites provados, não ótimo algébrico.

\paragraph{Regressão SoCG (48 casos, 4 por estrato)} O tempo total caiu de 131,1 s (fluxo exato, 08/10) para 23,5 s com o estágio \texttt{binary64} do oráculo (26,6 s na rodada de 09/10), 3,5 s com a triagem em \texttt{binary64} e 2,7 s com as provas pelo erro \emph{a priori}. Nas 134.300 chamadas do padrão de 08/10, nenhuma usou aritmética racional: 120.734 fecharam pela prova de um candidato da construção e 13.566 pelo polimento. Num replay de 11.411 chamadas, a mediana por chamada foi de 12–20~µs, contra 560–910~µs no fluxo exato.

\paragraph{Coleção da Paula (235 instâncias)} Em 60 s, sem tour inicial, fecharam 197 com o fluxo exato, 207 com o estágio \texttt{binary64} e 209 com a triagem em \texttt{binary64}. O pipeline de 600 s (ILS e B\&B preguiçoso sem mergulhos) fechou 17 das 38 abertas na versão de 07/10, e a campanha longa na armijo (\cref{sec:memoria}) fechou mais 6: somando todas as rodadas, 220 ótimos certificados. Nas 15 restantes (75 a 100 regiões), o gap provado vai de 0,70\% a 18,5\% (mediana 7,7\%), em parte medido com fronteiras de 1,5 GiB e, portanto, limitado pela memória.

\paragraph{Contra as outras abordagens} Na mesma coleção, o solver de Fekete et al.\ fechou 124 instâncias contra as nossas 197 (06/10), e onde ambos fecham fomos 7,0 vezes mais rápidos em média geométrica. O CPLEX do rascunho, com 1 h, provou 80 das 87 instâncias de até 15 regiões; nas 80, o nosso valor difere do dele no máximo \(3{,}2\times10^{-5}\), e as 7 que ele deixou abertas nós provamos em menos de 0,05 s. Contra o ILS-BCD da Paula (melhor de 10 execuções, outra máquina): em todas as instâncias que certificamos, o melhor tour dela é ótimo, a menos do arredondamento em duas casas, e nas 15 abertas o nosso melhor tour empata com o dela. Em 100i1000-410, as 10 sementes da nossa ILS pararam 0,19\% acima; foi o B\&B que achou o tour dela e provou que é ótimo.

%==========================================================================
\section{Limitações e questões em aberto}\label{sec:limites}
\begin{itemize}[leftmargin=*,itemsep=2pt]
\item \textbf{Instâncias grandes.} Acima de cerca de 60 regiões, nem nós, nem Fekete, nem o modelo do CPLEX fecham as instâncias da Paula. Ramificação forte, \emph{lookahead}, DFS/BFS, partidas primais e LNS exata não fecharam mais; relaxações em grafo (Held--Karp com distâncias entre regiões, custos por triplas) são mais fracas que as de inserção mesmo na ordem ótima.
\item \textbf{O piso da fronteira.} O limite inferior congela quando o mínimo da fronteira alcança o piso dos nós descartados (\cref{fig:piso}), e a execução continua sem poder subi-lo. Desde 12/10, a execução termina quando isso acontece, e o nó custa menos (\cref{sec:memoria}). A correção estrutural continua em aberto: não perder os nós descartados, gravando-os em disco ou reexplorando as suas regiões por busca em profundidade limitada pelo piso.
\item \textbf{A ILS.} Não tem garantia. Em 100i1000-410, as 10 sementes de 1.200 s param 0,19\% acima do ótimo, que difere do tour delas sobretudo por um trecho de \(\approx\)16 regiões em outra posição; blocos maiores no Or-opt não ajudaram. Faltam vizinhanças dela (3-opt, Balas--Simonetti) e uma comparação na mesma máquina, com o mesmo orçamento.
\item \textbf{Exatidão.} ``Ótimo certificado'' quer dizer gap \(10^{-6}\) fechado com limites provados das relaxações e verificação de visitas com tolerância; não é um ótimo racional exato do problema não convexo. A prova geral da construção com interseção continua em aberto~\cite{tutorial}; os limites não dependem dela.
\item \textbf{Medições.} Uma máquina, uma repetição por variante, várias execuções simultâneas em algumas rodadas; nos casos de milissegundos, as razões são ruído.
\end{itemize}

%==========================================================================
\appendix
\section{Notação}
\begin{center}\small
\begin{tabular}{@{}ll@{}}\toprule
\(R_1,\dots,R_n\) & regiões da instância (pontos, segmentos, polígonos simples)\\
\(P_1,\dots,P_k\) & regiões convexas de uma relaxação, na ordem cíclica\\
\(p_i,q_i\) & contato com \(P_i\); índices cíclicos, \(p_{k+1}=p_1\)\\
\(d_i,u_i\) & elo \(p_{i+1}-p_i\) e o seu vetor dual, \(\norm{u_i}\le1\)\\
\(c_i\) & coeficiente \(u_{i-1}-u_i\) do contato \(i\)\\
\(\supp_i(c)\) & suporte \(\min_{v\in P_i}(v-r)\cdot c\)\\
\(D(u)\) & limite dual do ciclo \eqref{eq:dualciclo}\\
\(F(x)\), \(f(t)\) & função âncora e a sua restrição a um trecho de aresta\\
\(\tau,c,g\) & gap e corte de uma chamada; gap do nó\\
\(U^\star\) & comprimento do incumbente\\
\(D\) (fronteira) & menor limite descartado, piso do limite inferior global\\
\(\mu,F\) & parâmetro de barreira; número total de faces\\
\(\eta\) & limiar de aceitação da ILS\\\bottomrule
\end{tabular}
\end{center}

\section{Guia do código}
\begin{center}\small
\begin{tabular}{@{}>{\raggedright\arraybackslash}p{42mm}>{\raggedright\arraybackslash}p{103mm}@{}}\toprule
Peça & Onde\\\midrule
Entrada e âncora de ponto & \code{tpp\_nonconvex\_tspn\_solve} em \path{packages/nonconvex-tpp/cpp/src/solvers/unordered.cpp}\\
Busca (as duas topologias) & \code{solve\_normalized\_unordered\_tpp}, \code{separated\_cycle\_root} e a classe \code{Frontier} (orçamento de memória), no mesmo arquivo\\
Adaptador do oráculo & \code{solve\_relaxation}, no mesmo arquivo\\
ILS & bloco \code{primal\_ils} de \code{solve\_normalized\_unordered\_tpp}; \code{best\_contact} em \path{unordered_geometry.cpp}\\
Triagem e inserções múltiplas & \code{insertion\_lower\_bounds}, \code{cycle\_multi\_insertion\_bound} em \path{unordered_bounds.cpp}\\
Dual em \texttt{binary64} & \code{tpp\_convex\_binary\_cycle\_dual}, \code{tpp\_convex\_binary\_insertion\_gain}, \code{tpp\_convex\_binary\_cycle\_insertion\_bounds} em \path{packages/convex-tpp/cpp/src/solvers/binary_dual.cpp}; suportes em \path{binary_dual.h}\\
Oráculo de ciclo & \code{tpp\_convex\_solve\_cycle\_double} e \code{tpp\_convex\_solve\_cycle} em \path{cycle.cpp}; construção em \path{cycle_refinement.h}\\
Estágio \texttt{binary64} & \code{tpp\_convex\_solve\_cycle\_float\_certified} em \path{cycle_float_oracle.cpp}; provas em \path{float_proof.h}; polimento em \path{float_chain.h}\\
Certificado exato & \code{tpp\_convex\_verify\_cycle\_certificate} em \path{cycle_certificate.cpp}; blocos de elos zero em \path{zero_contact_certificate.h}\\
Campanha da Paula & \path{benchmarks/_internal/paula_campaign.py} (\code{tpp.py paula-campaign})\\\bottomrule
\end{tabular}
\end{center}

\clearpage
\begin{thebibliography}{9}\small\raggedright
\bibitem{dror} M.~Dror, A.~Efrat, A.~Lubiw e J.~S.~B. Mitchell. \emph{Touring a sequence of polygons}. Proceedings of STOC 2003, pp.~473--482. \href{https://doi.org/10.1145/780542.780612}{doi:10.1145/780542.780612}.
\bibitem{fekete} S.~P. Fekete, R.~Kniep, D.~Krupke e M.~Perk. \emph{A branch-and-bound algorithm for the traveling salesman problem with difficult neighborhoods}. SoCG 2026 (versão anterior no EuroCG 2025). Solver usado: fork em \path{third_party/tspn-socg}.
\bibitem{paula} Rascunho do artigo sobre o TSPN com vizinhanças poligonais não convexas (ILS-BCD e modelo MIQCP), versão de 07/10/2026. Material local, sem autorização de redistribuição.
\bibitem{oraculo} TouringPolygons. \emph{O oráculo convexo certificado}, relatório técnico, 9 de outubro de 2026. \path{docs/reports/oraculo-convexo/main.tex}.
\bibitem{ordemlivre} TouringPolygons. \emph{TPP não convexo com ordem livre}. \path{docs/algorithms/unordered-tpp.md}; o \cref{lem:apriori} está na seção ``Certificado convexo e interseções''.
\bibitem{ciclo} TouringPolygons. \emph{Exact ordered cycles through closed convex polygons}. \path{docs/algorithms/convex-cycle.md}.
\bibitem{certificadociclo} TouringPolygons. \emph{Exact certificate for an ordered convex cycle}. \path{docs/algorithms/convex-cycle-certificate.md}.
\bibitem{tspnmd} TouringPolygons. \emph{TSPN with the maintained insertion/decomposition branch and bound}. \path{docs/algorithms/tspn.md}.
\bibitem{tutorial} TouringPolygons. \emph{A tutorial on fixed-order convex touring polygons and directional last-step maps}, 2026. \path{docs/reports/directional-tpp-tutorial/main.tex}.
\bibitem{contraexemplos} TouringPolygons. \emph{Unordered pruning counterexamples}. \path{docs/algorithms/unordered-pruning-counterexamples.md}.
\bibitem{resultados} TouringPolygons. Seções \code{tspn-paula-cycle-2026-10-06}, \code{tspn-paula-lower-bound-2026-10-07}, \code{tspn-cycle-float-2026-10-08}, \code{tspn-cycle-insertion-binary64-2026-10-09} e \code{tspn-search-retune-2026-10-09} de \path{benchmarks/results-saved/README.md}.
\end{thebibliography}
\end{document}
